\section{表示、原子逼近与进入泛函分析}\label{sec:07-04}

\subsection{Daniell--Stone 与 Riesz--Markov：从正泛函到 Radon 测度}\label{subsec:7-4-1}
\index{Daniell--Stone 定理}\index{Riesz--Markov 表示定理}\index{Radon 测度!正泛函表示}

测度通常先在集合上给出，再由简单函数逼近定义积分。Daniell 的路线反过来：先规定一族可观测函数的
“积分”，再追问这些数值是否来自一个可数可加测度。有限可加并不困难；真正把有限信息升级为测度的，
是单调列下降到零时积分也下降到零。Riesz--Markov 定理则把拓扑加入这条路线：紧支连续函数不仅生成
Borel 信息，还能借助局部 Urysohn 引理把紧集和开集夹在同一个测试函数中，因而恢复 Radon 正则性。

\begin{definition}[Daniell 泛函]\label{def:7-4-1-1}
设 $X$ 为集合，$\mathcal L$ 是 $X$ 上含常数的有界实函数向量格，即
$f,g\in\mathcal L$ 蕴含 $af+bg,f\vee g,f\wedge g\in\mathcal L$。正线性泛函
$I:\mathcal L\to\R$ 若满足
\[
 f_n\downarrow0\quad\Longrightarrow\quad I(f_n)\downarrow0,
\]
就称为 \emph{Daniell 泛函}。这里的下降和极限均为逐点意义；这个条件称为 Daniell 连续性。
\end{definition}

向量格条件还蕴含 $|f|,f^+,f^-$ 以及对任意连续分段线性函数 $q:\R\to\R$ 的复合
$q\circ f$（只需在 $f(X)$ 所在的有界区间上分段线性）仍属于 $\mathcal L$。这些截断正是外测度
构造中切分覆盖的工具。

\begin{definition}[正 Radon 泛函]\label{def:7-4-1-2}
设 $X$ 为局部紧 Hausdorff 空间。线性泛函 $\Lambda:C_c(X;\R)\to\R$ 若满足
\[
 f\ge0\quad\Longrightarrow\quad\Lambda(f)\ge0,
\]
就称为正 Radon 泛函。它不要求 $C_c(X)$ 带一个全局 Banach 范数；但若 $K$ 紧，取
$\chi\in C_c(X,[0,1])$ 在 $K$ 上等于 $1$，则对 $\supp f\subset K$ 有
\begin{equation}\label{eq:7-4-1-local-bound}
 |\Lambda(f)|\le \Lambda(\chi)\norm f_\infty.
\end{equation}
\end{definition}

确实，$-\norm f_\infty\chi\le f\le\norm f_\infty\chi$，分别对两边应用正性便得到
\eqref{eq:7-4-1-local-bound}。所以正性自动给出每个固定紧支撑层上的一致范数有界性，而没有虚构
$C_c(X)$ 的单一全局上确界范数结构。

\begin{definition}[表示测度]\label{def:7-4-1-3}
若 Borel 测度 $\mu$ 满足
\[
 \Lambda(f)=\int_X f\dd\mu\qquad(f\in C_c(X)),
\]
则称 $\mu$ 表示 $\Lambda$。本节所称 Radon 测度均满足：紧集有限、开集内正则、Borel 集外正则。
\end{definition}

\begin{theorem}[Daniell--Stone]\label{theorem:7-4-1-1}
设 $\mathcal L$ 是集合 $X$ 上含常数的有界实函数向量格，$I:\mathcal L\to\R$ 正线性，
$I(1)<\infty$，并具有 Daniell 连续性。则在 $\sigma(\mathcal L)$ 上存在唯一有限测度 $\mu$，使
\[
 I(f)=\int_X f\dd\mu\qquad(f\in\mathcal L).
\]
\end{theorem}

\begin{proof}
\textbf{第一步：可数函数覆盖成本。}
对任意函数 $h:X\to[0,\infty]$ 定义 Daniell 上积分
\begin{equation}\label{eq:7-4-1-daniell-upper-integral}
 J(h)=\inf\left\{
   \sum_{n=1}^\infty I(g_n):
   g_n\in\mathcal L_+,\quad
   h\le\sum_{n=1}^\infty g_n
 \right\}.
\end{equation}
若 $g\in\mathcal L_+$ 且 $g\le\sum_ng_n$，置
\[
 r_N=\left(g-\sum_{n=1}^Ng_n\right)^+.
\]
则 $r_N\in\mathcal L_+$ 且 $r_N\downarrow0$。又有
$g\le\sum_{n=1}^Ng_n+r_N$，所以正性与 Daniell 连续性给
\begin{equation}\label{eq:7-4-1-cover-dominates}
 I(g)\le\sum_{n=1}^NI(g_n)+I(r_N)
 \longrightarrow\sum_{n=1}^\infty I(g_n).
\end{equation}
对所有覆盖取下确界得到 $I(g)\le J(g)$；以 $g$ 自身作单项覆盖又给 $J(g)\le I(g)$，故
\begin{equation}\label{eq:7-4-1-J-agrees-I}
 J(g)=I(g)\qquad(g\in\mathcal L_+).
\end{equation}

现在定义
\begin{equation}\label{eq:7-4-1-daniell-outer}
 \mu^*(A)=J(\1_A)\qquad(A\subset X).
\end{equation}
空覆盖给 $\mu^*(\varnothing)=0$，单调性来自可用覆盖的包含关系。又因常数函数 $1$ 覆盖每个
$A\subset X$，有 $\mu^*(A)\leq I(1)<\infty$。若
$A\subset\bigcup_kA_k$，为每个 $A_k$ 取成本在
$\mu^*(A_k)+\varepsilon2^{-k}$ 内的可数函数覆盖，再把双重列枚举成一列，便得到
\[
 \mu^*(A)\le\sum_k\mu^*(A_k)+\varepsilon.
\]
令 $\varepsilon\downarrow0$，得到可数次可加性。因此 $\mu^*$ 是外测度；并且由
\eqref{eq:7-4-1-J-agrees-I}，
$\mu^*(X)=J(1)=I(1)<\infty$。这一步说明为什么不能用单个主控函数的成本下确界代替
\eqref{eq:7-4-1-daniell-upper-integral}：可数次可加性来自把各集合的可数覆盖合并，而不是来自一次
有限格运算。

\textbf{第二步：半直线逆像满足 \Caratheodory{} 条件。}
固定 $f\in\mathcal L$、$c\in\R$，令 $E=\{f>c\}$。取任意 $B\subset X$，并取一个有限成本覆盖
\[
 \1_B\le\sum_{n=1}^\infty g_n,\qquad
 S=\sum_{n=1}^\infty I(g_n)<\infty.
\]
这种覆盖总能取到，因为常数 $1$ 覆盖 $X$。令
\[
 u_k=\bigl(k(f-c)^+\bigr)\wedge1,\qquad
 M_n=\norm{g_n}_\infty,\qquad
 a_{n,k}=g_n\wedge(M_n\,u_k),\qquad a_{n,0}=0.
\]
则 $a_{n,k}\uparrow g_n\1_E$ 逐点成立。增量
$d_{n,k}=a_{n,k}-a_{n,k-1}$ 属于 $\mathcal L_+$，而双重族 $(d_{n,k})_{n,k}$ 覆盖 $B\cap E$，因为
\[
 \sum_{n,k}d_{n,k}(x)=\sum_ng_n(x)\quad(x\in E).
\]
记
\[
 D=\sum_n\lim_{k\to\infty}I(a_{n,k}).
\]
由于 $0\le a_{n,k}\le g_n$，有 $0\le D\le S$；非负级数的单调收敛还给
\[
 \mu^*(B\cap E)\le\sum_{n,k}I(d_{n,k})=D.
\]
另一方面，对固定 $k$ 令 $b_{n,k}=g_n-a_{n,k}$。在 $E^c$ 上 $u_k=0$，故 $(b_{n,k})_n$ 覆盖
$B\setminus E$，于是
\[
 \mu^*(B\setminus E)
 \le\sum_nI(b_{n,k})
 =S-\sum_nI(a_{n,k}).
\]
令 $k\to\infty$，非负级数的单调收敛给
\[
 \mu^*(B\cap E)+\mu^*(B\setminus E)\le D+(S-D)=S.
\]
最后对 $B$ 的覆盖成本取下确界，得到 \Caratheodory{} 反向不等式；正向由外测度次可加性成立。
所以每个 $\{f>c\}$ 都可测。这样的半直线逆像生成 $\sigma(\mathcal L)$，故
$\sigma(\mathcal L)$ 包含在 \Caratheodory{} 可测类中。由定理~\ref{theorem:2-1-2-2}，
$\mu=\mu^*|_{\sigma(\mathcal L)}$ 是有限测度。

\textbf{第三步：层级逼近恢复 $I$。}
仍令 $E_t=\{f>t\}$，其中 $f\in\mathcal L$。函数
\[
 u_{t,k}=\bigl(k(f-t)^+\bigr)\wedge1
\]
递增到 $\1_{E_t}$，并约定 $u_{t,0}=0$。由 \eqref{eq:7-4-1-J-agrees-I} 和单调性，
$I(u_{t,k})=J(u_{t,k})\le J(\1_{E_t})=\mu(E_t)$。反向把增量
$u_{t,k}-u_{t,k-1}$ 枚举成一个可数 $\mathcal L_+$ 覆盖；其和在 $E_t$ 上等于 $1$，故
\begin{equation}\label{eq:7-4-1-level-functional}
 \mu(E_t)\le\sum_{k=1}^\infty I(u_{t,k}-u_{t,k-1})
 =\lim_{k\to\infty}I(u_{t,k})\le\mu(E_t).
\end{equation}
为完整写出闭水平集一侧，令
\[
 w_{t,k}=\bigl(k(t-f)^+\bigr)\wedge1.
\]
把刚证的开水平集公式用于函数 $-f$ 和阈值 $-t$，得到
$I(w_{t,k})\uparrow\mu(\{f<t\})$。因此，对
\[
 v_{t,k}=1-w_{t,k}
\]
有
\begin{equation}\label{eq:7-4-1-closed-level-functional}
 I(v_{t,k})
 =I(1)-I(w_{t,k})
 \downarrow\mu(X)-\mu(\{f<t\})
 =\mu(\{f\ge t\}).
\end{equation}

现在设 $0\le f\le M$。给定 $\delta>0$，只把求和写到 $j\delta\le M+\delta$，并置
\[
 s_\delta=\delta\sum_{j\ge1}\1_{\{f>j\delta\}},
 \qquad
 t_\delta=\delta\1_X+
 \delta\sum_{j\ge1}\1_{\{f\ge j\delta\}}.
\]
这是有限和。固定 $x$ 并记 $y=f(x)/\delta$，则
\[
 \#\{j\geq1:j<y\}\leq y
 \leq1+\#\{j\geq1:j<y\},
\]
以及
\[
 y\leq1+\#\{j\geq1:j\leq y\}\leq y+1.
\]
分别乘以 $\delta$，便逐点得到
\[
 0\le f-s_\delta\le\delta,\qquad
 0\le t_\delta-f\le\delta.
\]
对下方每个水平使用 $u_{j\delta,k}$，有限求和后有
$\delta\sum_ju_{j\delta,k}\le s_\delta\le f$。令 $k\to\infty$，由
\eqref{eq:7-4-1-level-functional} 得
\[
 \int_Xs_\delta\dd\mu
 =\delta\sum_j\mu(f>j\delta)\le I(f).
\]
对上方水平使用 $v_{j\delta,k}$，有
$f\le\delta1+\delta\sum_jv_{j\delta,k}$；令 $k\to\infty$ 并用
\eqref{eq:7-4-1-closed-level-functional}，得到
\[
 I(f)\le\delta I(1)+\delta\sum_j\mu(f\ge j\delta)
 =\int_Xt_\delta\dd\mu.
\]
因为 $\mu(X)=I(1)<\infty$，上下阶梯与 $f$ 的一致误差至多 $\delta$。令 $\delta\downarrow0$，夹逼给
$I(f)=\int f\dd\mu$。一般实值 $f$ 写成 $f^+-f^-$ 即得表示。

\textbf{第四步：唯一性。}
若有限测度 $\nu$ 也表示 $I$，则对 $A\in\sigma(\mathcal L)$ 的任意函数覆盖
$\1_A\le\sum_ng_n$，积分的单调收敛性给
\[
 \nu(A)\le\sum_n\int g_n\dd\nu=\sum_nI(g_n).
\]
对覆盖取下确界，得 $\nu(A)\le\mu^*(A)=\mu(A)$。又
$\nu(X)=I(1)=\mu(X)$；把同一不等式用于 $A^c$，便得到反向不等式，故 $\nu(A)=\mu(A)$。
\end{proof}

若函数格不含常数或 $I(1)=\infty$，应先在可由某个 $h\in\mathcal L_+$ 主控的局部层上构造，再核对
这些局部测度的相容性。不能把有限版定理中的 $1$ 直接删去而保留原证明。

下面进入拓扑版本。为避免把正则性藏在一句“标准构造”中，我们从开集上的测试函数上确界开始。

\begin{theorem}[Riesz--Markov 表示]\label{theorem:7-4-1-2}
设 $X$ 为局部紧 Hausdorff 空间。每个正线性泛函
$\Lambda:C_c(X;\R)\to\R$ 都由唯一 Radon 测度 $\mu$ 表示。反之，每个 Radon 测度都在
$C_c(X;\R)$ 上定义正线性泛函。
\end{theorem}

\begin{proof}
\textbf{第一步：开集上的函数上确界及其可加性。}
对开集 $U\subset X$ 定义
\begin{equation}\label{eq:7-4-1-open-mass}
 p(U)=\sup\{\Lambda(f):f\in C_c(X),\ 0\le f\le1,\ \supp f\subset U\}.
\end{equation}
零函数是 $U=\varnothing$ 时的唯一候选，故 $p(\varnothing)=0$；若 $U\subset V$，$U$ 的每个候选也是
$V$ 的候选，故 $p(U)\le p(V)$。若 $U$ 相对紧，取
$\chi\in C_c(X,[0,1])$ 在 $\overline U$ 上等于 $1$，则每个候选 $f$ 满足 $f\le\chi$，从而
$p(U)\le\Lambda(\chi)<\infty$。

若 $U\subset\bigcup_{j\ge1}U_j$ 且 $f$ 是 \eqref{eq:7-4-1-open-mass} 中的候选函数，紧集
$\supp f$ 被有限多个 $U_j$ 覆盖。命题~\ref{proposition:1-3-1-finite-partition-of-unity}
给出相应的有限单位分解 $(\varphi_j)$；于是
\[
 f=\sum_jf\varphi_j,\qquad
 0\le f\varphi_j\le1,\qquad
 \supp(f\varphi_j)\subset U_j.
\]
正性给 $\Lambda(f)\le\sum_jp(U_j)$，取上确界得到
\begin{equation}\label{eq:7-4-1-open-subadd}
 p(U)\le\sum_{j=1}^\infty p(U_j).
\end{equation}
若 $U,V$ 不交，分别取支撑在其中的候选函数 $f,g$，则 $0\le f+g\le1$ 且支撑在 $U\cup V$，故
$p(U\cup V)\ge p(U)+p(V)$。对两两不交的开集 $(U_j)$，归纳应用这项有限超可加性可得
\[
 p\left(\bigcup_{j\geq1}U_j\right)
 \geq p\left(\bigcup_{j=1}^NU_j\right)
 =\sum_{j=1}^Np(U_j)
\]
对每个 $N$ 成立；令 $N\to\infty$，再结合 \eqref{eq:7-4-1-open-subadd}，便得到可数可加性。

还需要开集的内部耗尽公式
\begin{equation}\label{eq:7-4-1-open-inner-p}
 p(U)=\sup\{p(V):V\text{ 开},\ \overline V\text{ 紧且 }\overline V\subset U\}.
\end{equation}
事实上，给定候选 $f$，局部紧性允许取开集 $V$ 使
$\supp f\subset V\Subset U$，于是 $\Lambda(f)\le p(V)$；反向由单调性得到。

\textbf{第二步：开外逼近产生外测度和 Borel 测度。}
对任意 $A\subset X$ 置
\begin{equation}\label{eq:7-4-1-riesz-outer}
 \mu^*(A)=\inf\{p(U):A\subset U,\ U\text{ 开}\}.
\end{equation}
空集和单调性直接成立。为证可数次可加性，若某个
$\mu^*(A_n)=\infty$，则所需不等式右端为无穷，结论直接成立。以下设每个
$\mu^*(A_n)<\infty$。给定 $\varepsilon>0$，由下确界定义可逐个选开集
$A_n\subset U_n$，使
$p(U_n)<\mu^*(A_n)+\varepsilon2^{-n}$；于是
\eqref{eq:7-4-1-open-subadd} 给
\[
 \mu^*\Bigl(\bigcup_nA_n\Bigr)
 \le p\Bigl(\bigcup_nU_n\Bigr)
 \le\sum_np(U_n)
 <\sum_n\mu^*(A_n)+\varepsilon.
\]
故 $\mu^*$ 是外测度，而且对开集 $U$ 有 $\mu^*(U)=p(U)$。

在检验 Borel 可测性之前，先建立紧集公式。若 $K$ 紧，则
\begin{equation}\label{eq:7-4-1-riesz-compact-majorant}
 \mu^*(K)=\inf\{\Lambda(\phi):
 \phi\in C_c(X),\ \phi\ge0,\ \phi|_K\ge1\}.
\end{equation}
记右端为 $q(K)$。给定开集 $U\supset K$，局部 Urysohn 引理给支撑在 $U$ 中且在 $K$ 上等于 $1$
的帽函数 $\phi$，所以 $q(K)\le\Lambda(\phi)\le p(U)$；对 $U$ 取下确界得到
$q(K)\le\mu^*(K)$。反向，若 $\phi$ 是右端候选且 $0<\eta<1$，开集
$U_\eta=\{\phi>1-\eta\}$ 包含 $K$。对每个支撑在 $U_\eta$ 中的 $0\le g\le1$，逐点有
$(1-\eta)g\le\phi$，故
\[
 \mu^*(K)\le p(U_\eta)\le\frac{\Lambda(\phi)}{1-\eta}.
\]
令 $\eta\downarrow0$，再对 $\phi$ 取下确界，得到另一向不等式。特别地，局部帽函数保证每个紧集的
外测度有限。

若紧集 $K,L$ 不交，取任意开集 $W\supset K\cup L$。在开集 $W\setminus L$ 中对紧集 $K$ 使用
局部 Urysohn 引理，取 $\phi\in C_c(X,[0,1])$，使 $\phi=1$ 于 $K$ 且
$\supp\phi\subset W\setminus L$。于是
\[
 U=W\cap\{\phi>2/3\},\qquad V=W\cap\{\phi<1/3\}
\]
是不交开集，分别包含 $K,L$，并且都包含于 $W$。集合函数 $p$ 在不交开集上的
超可加性给
\[
 p(W)\ge p(U)+p(V)\ge\mu^*(K)+\mu^*(L).
\]
对 $W$ 取下确界，再结合外测度次可加性，得到紧集有限可加性
\begin{equation}\label{eq:7-4-1-riesz-compact-additivity}
 \mu^*(K\cup L)=\mu^*(K)+\mu^*(L).
\end{equation}
此外，\eqref{eq:7-4-1-open-inner-p} 可升级成
\begin{equation}\label{eq:7-4-1-riesz-inner-regular}
 p(U)=\sup\{\mu^*(K):K\subset U,\ K\text{ 紧}\}.
\end{equation}
一向来自 $\mu^*(K)\le p(U)$。反向取 $V\Subset U$；对每个开集 $O\supset\overline V$ 有
$p(O)\ge p(V)$，故 $\mu^*(\overline V)\ge p(V)$，再用
\eqref{eq:7-4-1-open-inner-p} 对 $V$ 取上确界。

下面专门证明开集满足 \Caratheodory{} 条件。先由
\eqref{eq:7-4-1-riesz-compact-additivity} 和
\eqref{eq:7-4-1-riesz-inner-regular} 得到：若 $K\subset V$、$K$ 紧而 $V$ 开，则
\begin{equation}\label{eq:7-4-1-riesz-remove-compact}
 p(V)\ge\mu^*(K)+p(V\setminus K).
\end{equation}
事实上，对任意紧集 $L\subset V\setminus K$，
$\mu^*(K)+\mu^*(L)=\mu^*(K\cup L)\le p(V)$；再对 $L$ 取上确界即可。

固定开集 $G$ 和 $A\subset X$，只需处理 $\mu^*(A)<\infty$。取开集 $V\supset A$ 使
$p(V)<\mu^*(A)+\varepsilon$，并令 $W=V\cap G$。由
\eqref{eq:7-4-1-riesz-inner-regular} 取紧集 $K\subset W$，使
$\mu^*(K)>p(W)-\varepsilon$。因为 $A\cap G\subset W$ 且
$A\setminus G\subset V\setminus K$，由 \eqref{eq:7-4-1-riesz-remove-compact}，
\begin{align*}
 \mu^*(A\cap G)+\mu^*(A\setminus G)
 &\le p(W)+p(V\setminus K)\\
 &\le p(W)+p(V)-\mu^*(K)\\
 &<\mu^*(A)+2\varepsilon.
\end{align*}
令 $\varepsilon\downarrow0$ 得 \Caratheodory{} 反向不等式；正向是外测度次可加性。因此所有开集可测，
\Caratheodory{} 可测类包含 $\Borel(X)$。以下把限制在 Borel 集上的测度记为 $\mu$。

\textbf{第三步：Radon 正则性。}
外正则性就是 \eqref{eq:7-4-1-riesz-outer}；紧集有限性来自
\eqref{eq:7-4-1-riesz-compact-majorant}；开集内正则性来自
\eqref{eq:7-4-1-riesz-inner-regular}，因为对开集 $U$ 有 $\mu(U)=\mu^*(U)=p(U)$。所以 $\mu$ 为
Radon 测度。这里所需的拓扑分离仅是紧集与开集之间的局部 Urysohn 引理，并不要求局部紧
Hausdorff 空间本身全局正规。

\textbf{第四步：从层集恢复 $\Lambda$。}
取 $0\le f\in C_c(X)$，令 $M=\norm f_\infty$、$K=\supp f$。若 $M=0$，则 $f=0$，积分表示由线性性成立。
对任意整数 $N\geq1$，令 $\delta=M/N$。对 $j=1,\ldots,N$，开集 $U_j=\{f>j\delta\}$ 的定义给出
$g_j\in C_c(U_j,[0,1])$，使
$\Lambda(g_j)>\mu(U_j)-\varepsilon/N$。这里
$\overline{U_j}\subset K\cap\{f\geq j\delta\}$ 是紧集，故
$\mu(U_j)\leq\mu(K)<\infty$，所以上确界的这一近似选择确实有意义。由于一个点至多落入
$\lfloor f(x)/\delta\rfloor$ 个 $U_j$，有 $\delta\sum_jg_j\le f$，故
\begin{equation}\label{eq:7-4-1-riesz-lower}
 \Lambda(f)\ge\delta\sum_{j=1}^N\mu(f>j\delta)-\delta\varepsilon.
\end{equation}

再令
\[
 h_j=\delta^{-1}\bigl((f-(j-1)\delta)^+\wedge\delta\bigr),
 \qquad f=\delta\sum_{j=1}^Nh_j.
\]
每个 $h_j$ 满足 $0\le h_j\le1$，其支撑包含于紧集
$K_j=K\cap\{f\ge(j-1)\delta\}$。对任意开集 $O\supset K_j$，$h_j$ 是 $p(O)$ 的候选函数，
故 $\Lambda(h_j)\le p(O)$；对 $O$ 取下确界得
$\Lambda(h_j)\le\mu(K_j)$。于是
\begin{equation}\label{eq:7-4-1-riesz-upper}
 \Lambda(f)\le\delta\sum_{j=1}^N\mu(K\cap\{f\ge(j-1)\delta\}).
\end{equation}
令 $G(t)=\mu(f>t)$，并记它在 $[0,M]$ 上的右、左端点和为
\[
 R_N^-:=\delta\sum_{j=1}^NG(j\delta),
 \qquad
 R_N^+:=\delta\sum_{j=0}^{N-1}G(j\delta).
\]
$G$ 有界且递减，故 Riemann 可积，并且
$R_N^-,R_N^+\to\int_0^MG(t)\dd t$。另一方面，令
\[
 U_N:=\delta\sum_{j=1}^N\mu\bigl(K\cap\{f\ge(j-1)\delta\}\bigr).
\]
由于正水平集均包含于 $K$，逐项相减得到
\begin{align*}
 0\le U_N-R_N^+
 &=\delta\bigl(\mu(K)-\mu(f>0)\bigr)
   +\delta\sum_{j=1}^{N-1}\mu(f=j\delta)\\
 &\le2\delta\mu(K).
\end{align*}
最后一个不等式使用了两两不交的水平集
$\{f=j\delta\}$。因此 $U_N\to\int_0^MG(t)\dd t$。把
\eqref{eq:7-4-1-riesz-lower}--\eqref{eq:7-4-1-riesz-upper} 写成
\[
 R_N^- -\delta\varepsilon\le\Lambda(f)\le U_N
\]
并令 $N\to\infty$（即 $\delta=M/N\downarrow0$），层蛋糕公式给
\[
 \Lambda(f)=\int_0^M\mu(f>t)\dd t=\int_Xf\dd\mu.
\]
实值 $f$ 由 $f^+-f^-$ 得到表示。

反过来，若 $\mu$ 为 Radon 测度，$f\in C_c(X)$ 的支撑紧，故
$\int|f|\dd\mu\le\norm f_\infty\mu(\supp f)<\infty$；积分的线性与正性给出正 Radon 泛函。

最后在本定理内完成唯一性。若另一个 Radon 测度 $\nu$ 也表示 $\Lambda$，则对开集 $U$，每个
$0\leq f\leq1$、$f\in C_c(U)$ 都满足 $\Lambda(f)\leq\nu(U)$，故 $p(U)\leq\nu(U)$。反过来，
对任意紧集 $K\subset U$ 取 $f\in C_c(U,[0,1])$ 且 $f=1$ 于 $K$，便有
$\nu(K)\leq\Lambda(f)\leq p(U)$；对 $K$ 使用开集内正则性得 $\nu(U)\leq p(U)$。因此
$\nu(U)=p(U)=\mu(U)$ 对所有开集成立。两测度的 Borel 外正则性随后给
$\nu(E)=\mu(E)$ 对每个 Borel 集 $E$ 成立。存在性、反向与唯一性均已证明；下一命题把刚用到的开集
恢复公式连同支撑刻画单独记录下来。
\end{proof}

\begin{proposition}[正则性与唯一性]\label{proposition:7-4-1-3}
若 $\mu$ 表示正泛函 $\Lambda$，则对每个开集 $U$，
\begin{equation}\label{eq:7-4-1-open-recovery}
 \mu(U)=\sup\{\Lambda(f):f\in C_c(U),\ 0\le f\le1\}.
\end{equation}
因此表示测度唯一。并且
\begin{equation}\label{eq:7-4-1-support-functional}
 x\notin\supp\mu
 \quad\Longleftrightarrow\quad
 \text{存在开邻域 }U\ni x\text{，使 }\Lambda(f)=0
 \text{ 对所有 }f\in C_c(U).
\end{equation}
\end{proposition}

\begin{proof}
若 $f$ 是 \eqref{eq:7-4-1-open-recovery} 中的候选函数，则
$\Lambda(f)=\int f\dd\mu\le\mu(U)$。反过来，给定紧集 $K\subset U$，局部 Urysohn 引理
\ref{theorem:1-3-1-2} 给 $f\in C_c(U,[0,1])$ 且 $f=1$ 于 $K$，所以
$\Lambda(f)\ge\mu(K)$。对 $K$ 取上确界并用开集内正则性，得到等号。

若 $\mu,\nu$ 都表示 $\Lambda$，开集公式给 $\mu(U)=\nu(U)$。对任意 Borel 集 $E$，外正则性给
\[
 \mu(E)=\inf_{U\supset E}\mu(U)
 =\inf_{U\supset E}\nu(U)=\nu(E),
\]
所以表示唯一。最后，若 $x\notin\supp\mu$，按支撑定义有开邻域 $U$ 满足 $\mu(U)=0$，支撑在
$U$ 中的函数积分均为零。反向若右侧成立，尤其所有非负候选函数的值为零；
\eqref{eq:7-4-1-open-recovery} 给 $\mu(U)=0$，故 $x\notin\supp\mu$。
\end{proof}

定理~\ref{theorem:7-3-1-2} 由 \Fejer{} 核和测度紧性证明 Herglotz 表示。现在给出另一种
正泛函证明。它所需的代数工具是 \Fejer{}--Riesz 分解：非负三角多项式可表为一个多项式的模平方。

\begin{lemma}[\Fejer{}--Riesz 分解]\label{lemma:7-4-1-fejer-riesz}
若三角多项式
\[
 q(t)=\sum_{k=-n}^na_ke^{ikt}
\]
在 $\T$ 上取实值且非负，则存在普通多项式 $P(z)=\sum_{k=0}^nb_kz^k$，使
\[
 q(t)=|P(e^{it})|^2.
\]
\end{lemma}

\begin{proof}
若 $q(t)\equiv a_0$，则 $a_0\ge0$；取常数多项式 $P\equiv\sqrt{a_0}$ 即得结论（包括
$a_0=0$）。以下设 $q$ 非常数，并删去两端的零系数，使 $a_n\ne0$；实值性给
$a_{-k}=\overline{a_k}$。令
\[
 Q(z)=z^n\sum_{k=-n}^na_kz^k.
\]
则 $Q$ 是次数 $2n$ 且常数项非零的多项式，并满足
\[
 Q(z)=z^{2n}\overline{Q(1/\overline z)}\qquad(z\ne0).
\]
所以不在单位圆上的零点按 $\alpha\leftrightarrow1/\overline\alpha$ 成对且重数相同。若
$e^{it_0}$ 是单位圆上的零点，则
$q(t)=e^{-int}Q(e^{it})$ 在 $t_0$ 的零点重数相同。设该重数为 $m$。由于 $q$ 是有限个指数函数的
线性组合，其 Taylor 展开在 $t_0$ 邻域成立；存在实数 $c\ne0$，使
\[
 q(t_0+h)=ch^m+o(|h|^m)\qquad(h\to0).
\]
若 $m$ 为奇数，右侧在 $h>0$ 与 $h<0$ 时异号，与 $q\geq0$ 矛盾；故 $m$ 为偶数（并且
$c>0$）。因此每个单位圆零点的重数均为偶数。

从每一对离圆零点中取一个，并从每个圆上零点取其一半重数，所得重数总和为 $n$。令 $P_0$ 以这些
零点为根，并置
$P_0^*(z)=z^n\overline{P_0(1/\overline z)}$；右式按 $P_0$ 的系数展开后在 $z=0$ 也延拓为
次数不超过 $n$ 的多项式。多项式 $Q$ 与
$P_0P_0^*$ 有完全相同的零点及重数，故 $Q=\lambda P_0P_0^*$。在 $|z|=1$ 上，
$P_0^*(z)=z^n\overline{P_0(z)}$，从而
$q(t)=\lambda|P_0(e^{it})|^2$。由 $q\ge0$ 且 $q\not\equiv0$ 得 $\lambda>0$；取
$P=\sqrt\lambda P_0$ 即可。
\end{proof}

\begin{proposition}[Herglotz 表示的正泛函证明]\label{proposition:7-4-1-herglotz-functional}
若复数列 $(c_n)_{n\in\Z}$ 满足 $c_{-n}=\overline{c_n}$，且每个 Toeplitz 矩阵
$T_N=(c_{j-k})_{0\le j,k\le N}$ 都半正定，则存在唯一的有限正 Borel 测度 $\mu$ 于 $\T$，使
\[
 c_n=\int_\T e^{-int}\dd\mu(t)\qquad(n\in\Z).
\]
\end{proposition}

\begin{proof}
在三角多项式代数 $\mathcal P$ 上定义
\[
 L\left(\sum_{k=-N}^Na_ke^{ikt}\right)=\sum_{k=-N}^Na_kc_{-k}.
\]
若 $r(t)=\sum_{k=0}^Nb_ke^{ikt}$，则
\[
 L(|r|^2)=\sum_{j,k=0}^N\overline{b_j}c_{j-k}b_k\ge0.
\]
由引理~\ref{lemma:7-4-1-fejer-riesz}，每个非负三角多项式都是某个 $|r|^2$，所以 $L$ 为正泛函。
关系 $c_{-n}=\overline{c_n}$ 先给出
$L(\overline h)=\overline{L(h)}$（$h\in\mathcal P$）。再把所用的 Cauchy--Schwarz 不等式写明。
给定 $p,q\in\mathcal P$，记
\[
 a=L(|p|^2),\qquad b=L(|q|^2),\qquad c=L(\overline q p).
\]
由 $L(|p+\lambda q|^2)\geq0$，对每个 $\lambda\in\C$ 有
\[
 0\leq a+\lambda\overline c+\overline\lambda c+|\lambda|^2b.
\]
若 $b>0$，取 $\lambda=-c/b$ 得 $|c|^2\leq ab$；若 $b=0$，先取
$\lambda=t\in\R$，再取 $\lambda=it$，并令 $t$ 在 $\R$ 中任意变化，分别迫使
$\Re c=0$ 与 $\Im c=0$。所以零分母情形也有 $c=0$，从而总有
\[
 |L(\overline q p)|^2\leq L(|p|^2)L(|q|^2).
\]
取 $q=1$，再由 $\norm p_\infty^2-|p|^2\ge0$，便有
\[
 |L(p)|^2\le L(1)L(|p|^2)
 \le c_0^2\norm p_\infty^2.
\]
故 $L$ 关于一致范数有界。由推论~\ref{corollary:7-2-1-2}，$\mathcal P$ 在 $C(\T)$ 中稠密，
于是 $L$ 唯一延拓为 $C(\T)$ 上的有界线性泛函 $\widetilde L$。该延拓仍为正：若
$f\in C(\T)$ 且 $f\ge0$，取 $p_m\in\mathcal P$ 一致逼近 $\sqrt f$，则
$|p_m|^2\to f$ 一致且 $L(|p_m|^2)\ge0$。

正性保证 $\widetilde L$ 在实值函数上取实值；又因 $\T$ 紧，
$C_c(\T;\R)=C(\T;\R)$。把 $\widetilde L$ 限制到这个实空间，并应用
定理~\ref{theorem:7-4-1-2}，得到唯一有限正测度 $\mu$。再由复线性性，
$\widetilde L(f)=\int f\dd\mu$ 对复值 $f\in C(\T)$ 也成立。特别地
\[
 \int_\T e^{-int}\dd\mu(t)=L(e^{-int})=c_n.
\]
若另一有限正测度具有相同的 Fourier 矩，则其积分泛函与 $L$ 在 $\mathcal P$ 上相同；三角多项式
稠密性和两个积分泛函的一致范数连续性说明它也表示 $\widetilde L$，故 Riesz--Markov 的唯一性迫使
它等于 $\mu$。
因此 Toeplitz 正性先在三角多项式上产生正泛函，一致范数连续性把它延到 $C(\T)$，
最后由 Riesz--Markov 定理得到表示测度。
\end{proof}

\begin{figure}[htbp]
\centering
\begin{tikzpicture}[node distance=7mm and 36mm]
  \node[book strong node] (fun) {正泛函\\$\Lambda:C_c(X)\to\R$};
  \node[book strong node,below=of fun] (open) {开集上确界\\$p(U)=\sup\Lambda(f)$};
  \node[book strong node,below=of open] (outer) {外测度\\$\mu^*(A)=\inf_{U\supset A}p(U)$};
  \node[book strong node,below=of outer] (radon) {Radon 测度\\内正则、外正则};
  \node[book warm node,below=of radon] (int) {积分表示\\$\Lambda(f)=\int f\dd\mu$};
  \node[book warm node,left=of outer] (ury) {Urysohn 帽函数\\$K\subset U$};
  \draw[book accent arrow] (fun)--(open);
  \draw[book accent arrow] (open)--(outer);
  \draw[book accent arrow] (outer)--(radon);
  \draw[book accent arrow] (radon)--(int);
  \draw[book dashed arrow,bend left=18] (ury.north east) to
    node[above,book label]{紧开夹逼}(open.west);
  \draw[book dashed arrow,bend right=18] (ury.south east) to
    node[below,book label]{层集阶梯和}(int.west);
\end{tikzpicture}
\caption{Riesz--Markov 构造链。开集公式产生外测度；Urysohn 函数同时负责局部有限、内正则与积分识别。}
\label{fig:7-4-1-riesz-route}
\end{figure}

图~\ref{fig:7-4-1-riesz-route} 中两条虚线是拓扑真正进入证明的位置。若只保留外测度箭头而删去
紧开夹逼，最多得到某个 Borel 测度，不能推出表示测度具有所需的 Radon 正则性。

\begin{theorem}[$C_0(X)^*$ 与复测度]\label{theorem:7-4-1-4}
设 $X$ 为局部紧 Hausdorff 空间，标量域为 $\R$ 或 $\C$。每个有界线性泛函
$\Lambda\in C_0(X)^*$ 唯一表示为
\[
 \Lambda(f)=\int_X f\dd\mu\qquad(f\in C_0(X)),
\]
其中 $\mu$ 是有限符号 Radon 测度（实情形）或有限复 Radon 测度（复情形），并且
\begin{equation}\label{eq:7-4-1-dual-variation-norm}
 \norm\Lambda=|\mu|(X).
\end{equation}
反之，每个这样的测度都定义一个有界线性泛函。
\end{theorem}

\begin{proof}
\textbf{第一步：先在泛函层作 Jordan 分解。}
先设标量为实数。对 $f\in C_0(X)_+$ 定义
\begin{equation}\label{eq:7-4-1-functional-positive-part}
 \Lambda^+(f)=\sup\{\Lambda(g):0\le g\le f\}.
\end{equation}
候选集非空，且
$|\Lambda(g)|\le\norm\Lambda\norm f_\infty$，所以上确界有限。正齐次性直接成立。若
$f_1,f_2\ge0$，候选 $g_i\le f_i$ 给
$\Lambda^+(f_1+f_2)\ge\Lambda^+(f_1)+\Lambda^+(f_2)$。反向取
$0\le h\le f_1+f_2$，令
\[
 h_1=h\wedge f_1,\qquad h_2=h-h_1.
\]
则 $0\le h_1\le f_1$，而
$h_2=(h-f_1)^+\le f_2$，故
$\Lambda(h)\le\Lambda^+(f_1)+\Lambda^+(f_2)$。对 $h$ 取上确界，得到可加性。

每个 $u\in C_0(X;\R)$ 都可写成两个正函数之差。若
$u=a-b=c-d$，其中 $a,b,c,d\ge0$，则 $a+d=b+c$；正锥上的可加性给
$\Lambda^+(a)-\Lambda^+(b)=\Lambda^+(c)-\Lambda^+(d)$。因此群完备化公式
\[
 \Lambda^+(a-b)=\Lambda^+(a)-\Lambda^+(b)
\]
良定义且给出实线性延拓；取 $a=u^+,b=u^-$ 就得到通常公式。若 $u\ge0$，
$0\le\Lambda^+(u)\le\norm\Lambda\norm u_\infty$，再对一般 $u$ 使用
$|\Lambda^+(u)|\le\Lambda^+(|u|)$，可知延拓有界且正。又因
\eqref{eq:7-4-1-functional-positive-part} 的候选中可取 $g=f$，
\[
 \Lambda^-(f):=\Lambda^+(f)-\Lambda(f)\ge0\qquad(f\ge0).
\]
所以 $\Lambda=\Lambda^+-\Lambda^-$ 先在泛函层分解为两个正泛函之差，随后可分别对
$\Lambda^+$ 和 $\Lambda^-$ 应用正 Riesz--Markov 定理。

为了把 $C_c(X)$ 上的表示延到 $C_0(X)$，先记录所需的稠密性。给 $u\in C_0(X)$ 和 $\varepsilon>0$，紧集
$K_\varepsilon=\{|u|\ge\varepsilon\}$ 有一个 $\chi\in C_c(X,[0,1])$ 满足
$\chi=1$ 于 $K_\varepsilon$；于是
$u\chi\in C_c(X)$ 且 $\norm{u-u\chi}_\infty<\varepsilon$。这正是
命题~\ref{proposition:1-3-1-4}。把
$\Lambda^\pm$ 限制到 $C_c(X)$，正 Riesz--Markov 定理给 Radon 测度 $\mu^\pm$。它们有限，因为
\[
 \mu^\pm(X)
 =\sup_{0\le f\le1,\ f\in C_c}\Lambda^\pm(f)
 \le\norm{\Lambda^\pm}<\infty.
\]
令 $\mu=\mu^+-\mu^-$。表示先在 $C_c$ 上成立；积分泛函与 $\Lambda$ 都关于上确界范数连续，故由
$C_c$ 稠密性延到全部 $C_0$。这里一般正泛函 $C_c(X)\to\R$ 的表示测度可以有无限总质量；正是
$C_0$ 泛函的全局有界性给出了上式中的 $\mu^\pm(X)<\infty$。

\textbf{第二步：复泛函由两个实泛函合成。}
设 $\Lambda$ 复线性。对实值 $f\in C_0(X;\R)$ 令
\[
 A(f)=\Re\Lambda(f),\qquad B(f)=\Im\Lambda(f).
\]
$A,B$ 是有界实线性泛函。由第一步，存在有限符号 Radon 测度 $\alpha,\beta$ 表示它们。
置 $\mu=\alpha+i\beta$。对实值 $u,v$，复线性给
\begin{align*}
 \Lambda(u+iv)
 &=\Lambda(u)+i\Lambda(v)\\
 &=A(u)-B(v)+i\bigl(B(u)+A(v)\bigr)\\
 &=\int_Xu\dd\alpha-\int_Xv\dd\beta
   +i\left(\int_Xu\dd\beta+\int_Xv\dd\alpha\right)\\
 &=\int_X(u+iv)\dd(\alpha+i\beta).
\end{align*}
故 $\mu$ 表示 $\Lambda$。这个顺序仍是“先分解泛函、再得到测度”。

\textbf{第三步：范数等于全变差。}
测度极分解给 $\dd\mu=h\dd|\mu|$，其中 $|h|=1$ $|\mu|$-a.e.。对
$\norm f_\infty\le1$，
\[
 |\Lambda(f)|\le\int|f|\dd|\mu|\le|\mu|(X),
\]
所以 $\norm\Lambda\le|\mu|(X)$。

反向不能把可测相位 $\overline h$ 直接当成 $C_0$ 测试。由
定理~\ref{theorem:3-3-3-1}，存在 $g_n\in C_c(X;\C)$ 使
$g_n\to\overline h$ 于 $L^1(|\mu|)$。令闭单位圆盘的径向投影为
\[
 \rho(z)=\begin{cases}z,&|z|\le1,\\ z/|z|,&|z|>1.
 \end{cases}
\]
则 $f_n=\rho\circ g_n\in C_c(X)$、$\norm{f_n}_\infty\le1$，并且因为径向投影固定单位圆周，
\[
 |f_n-\overline h|\le2|g_n-\overline h|.
\]
这个点态估计可分两种情况核对：$|g_n|\leq1$ 时 $f_n=g_n$，所需不等式立即成立；$|g_n|>1$ 时，由
$|\overline h|=1$ 与反三角不等式，
\[
 \left|\frac{g_n}{|g_n|}-\overline h\right|
 \leq\left|\frac{g_n}{|g_n|}-g_n\right|+|g_n-\overline h|
 =\bigl||g_n|-1\bigr|+|g_n-\overline h|
 \leq2|g_n-\overline h|.
\]
因此
\[
 \left|\Lambda(f_n)-|\mu|(X)\right|
 =\left|\int_X(f_nh-1)\dd|\mu|\right|
 \le\int_X|f_n-\overline h|\dd|\mu|\longrightarrow0.
\]
取上极限得到 $\norm\Lambda\ge|\mu|(X)$，与前一不等式合并即为
\eqref{eq:7-4-1-dual-variation-norm}。实符号情形可取 $h=\sgn\mu\in\{-1,1\}$，同一逼近完全适用。

若 $\mu,\nu$ 都表示同一泛函，则 $\mu-\nu$ 表示零泛函。把刚证明的范数公式
用于该差测度，得到 $|\mu-\nu|(X)=0$，所以 $\mu=\nu$。

最后，给定有限符号或复 Radon 测度，绝对值估计直接给
$|\int f\dd\mu|\le|\mu|(X)\norm f_\infty$，故它确实定义 $C_0(X)^*$ 中的泛函；把上面的
$L^1(|\mu|)$ 相位逼近原样应用于这个泛函，又给出反向范数不等式，所以其范数仍恰为 $|\mu|(X)$。
\end{proof}

\begin{example}[点值泛函]\label{ex:7-4-1-1}
固定 $x\in X$，令 $\Lambda_x(f)=f(x)$。对 Borel 集 $A$，Dirac 测度满足
\[
 \int_X f\dd\delta_x=f(x),\qquad
 \delta_x(A)=\1_A(x).
\]
故 $\delta_x$ 表示 $\Lambda_x$，且 $\norm{\Lambda_x}=1=|\delta_x|(X)$。更一般地，
$\sum_{j=1}^Nc_jf(x_j)$ 由 $\sum_jc_j\delta_{x_j}$ 表示；若 $c_j\ge0$，它就是正泛函。
\end{example}

\begin{example}[加权积分]\label{ex:7-4-1-2}
在 $X=\R^d$ 上设 $w\ge0$ 局部可积。对 $f\in C_c(\R^d)$，
\[
 \Lambda(f)=\int_{\R^d}f(x)w(x)\dd x
\]
有限且正，由 Radon 测度 $\mu=w\dd x$ 表示。更准确地，令
\[
 \operatorname{ess\,supp}w
 =\R^d\setminus\bigcup\{U:U\text{ 开且 }w=0\text{ a.e. 于 }U\}.
\]
这是一个闭集，并且
\[
 x\notin\supp\mu
 \Longleftrightarrow
 \text{存在开邻域 }U\ni x\text{ 使 }w=0\text{ a.e. 于 }U.
\]
所以 $\supp\mu=\operatorname{ess\,supp}w$。若另有 $w\in L^1$，则该泛函延到 $C_0$，且
$\norm\Lambda=\int|w|\dd x$；仅有局部可积时通常只在 $C_c$ 上定义。
\end{example}

\begin{example}[非正泛函]\label{ex:7-4-1-3}
取 $\phi\in C_c^1(\R)$ 且 $\phi'(0)\ne0$，令
$f_n(x)=n^{-1}\phi(nx)$。则
\[
 \norm{f_n}_\infty=n^{-1}\norm\phi_\infty\longrightarrow0,
 \qquad f_n'(0)=\phi'(0).
\]
因此 $f\mapsto f'(0)$ 在 $C_c^1$ 的
$\norm f_{C^1}=\norm f_\infty+\norm{f'}_\infty$ 拓扑下连续，却不可能连续延到带上确界范数的
$C_0(\R)$。若它由有限符号测度 $\nu$ 表示，就会有
$|f_n'(0)|\le|\nu|(\R)\norm{f_n}_\infty\to0$，与上式矛盾。测试空间的拓扑改变后，表示对象会从
测度升级为一阶分布。
\end{example}

两条表示路线的分工现在很清楚。Daniell--Stone 用单调连续性换取可数可加，适合从随机变量格或初等
积分出发；Riesz--Markov 用局部紧拓扑和紧支连续测试换取正则性，适合弱收敛和对偶空间。

\subsection{原子测度、经验测度与求积：有限对象逼近 Radon 测度}\label{subsec:7-4-2}
\index{原子测度!弱星稠密}\index{经验测度}\index{Tchakaloff 定理}\index{求积公式}

一个弱星基本邻域只询问有限多个积分。因此，为了进入这个邻域，并不需要逐点重建测度；只需把空间
分成有限片，使每个测试函数在每片上的振幅很小，再把整片质量搬到一个代表点。随机抽样把同一结构
变成经验测度，有限维凸性则把“近似匹配”加强为 Tchakaloff 的“精确匹配”。

\begin{example}[Riemann 求积]\label{ex:7-4-2-1}
设 $a=x_0<\cdots<x_N=b$，$\xi_j=(x_{j-1}+x_j)/2$，并令
\[
 \nu_P=\sum_{j=1}^N(x_j-x_{j-1})\delta_{\xi_j}.
\]
若 $|P|=\max_j(x_j-x_{j-1})$，连续函数 $f$ 的模连续性记为
$\omega_f(r)=\sup_{|x-y|\le r}|f(x)-f(y)|$，则
\begin{align*}
 \left|\int_a^bf(x)\dd x-\int f\dd\nu_P\right|
 &\le\sum_{j=1}^N\int_{x_{j-1}}^{x_j}|f(x)-f(\xi_j)|\dd x\\
 &\le(b-a)\omega_f(|P|/2)\longrightarrow0.
\end{align*}
这就是 Lebesgue 测度的有限原子逼近。若每段改用端点平均，得到梯形权重，完全相同的计算给误差
至多 $(b-a)\omega_f(|P|)$；Gauss 节点则不是只压低振幅误差，而是精确匹配一段有限维多项式空间。
\end{example}

\begin{example}[经验测度的测试积分]\label{ex:7-4-2-2}
若 $X_1,X_2,\ldots$ 独立同分布，共同分布为 $\mu$，则对任意有界实可测 $f$，
\[
 \int f\dd L_n=\frac1n\sum_{k=1}^nf(X_k),\qquad
 \E\int f\dd L_n=\int f\dd\mu,
\]
并且在 $f\in L^2(\mu)$ 时
\[
 \Var\!\left(\int f\dd L_n\right)=\frac{\Var(f(X_1))}{n}.
\]
单个 $f$ 的强大数律只给一个测试方向。Polish 假设的作用，是允许选取一个可数的
有界 Lipschitz 决定类，并对相应概率一事件取可数交；定理~\ref{theorem:7-4-2-3} 将完成这一步。
\end{example}

\begin{example}[矩匹配求积]\label{ex:7-4-2-3}
取 $K=[-1,1]$ 上的均匀概率 $\mu=\frac12\dd x$，并只观察
$V=\Span\{1,x,x^2\}$。两点测度
\[
 \nu=\frac12\delta_{-1/\sqrt3}+\frac12\delta_{1/\sqrt3}
\]
满足
\[
 \int1\dd\nu=1=\int1\dd\mu,\qquad
 \int x\dd\nu=0=\int x\dd\mu,\qquad
 \int x^2\dd\nu=\frac13=\int x^2\dd\mu.
\]
它没有逼近全部连续测试，却对指定的三维空间给出零误差。这正是 Gauss 两点求积的最低阶模型。
\end{example}

\begin{definition}[原子测度]\label{def:7-4-2-1}
形如
\[
 \nu=\sum_{j=1}^Nc_j\delta_{x_j}
\]
的符号或复测度称为有限原子测度。若 $c_j\ge0$ 且 $\sum_jc_j=1$，它是有限支持概率测度；对任意
可积测试函数，
$\int f\dd\nu=\sum_jc_jf(x_j)$。
\end{definition}

\begin{definition}[弱星测试邻域]\label{def:7-4-2-2}
把有限符号或复 Radon 测度依定理~\ref{theorem:7-4-1-4} 看作 $C_0(X)^*$ 的元素。
给定有限集 $G\subset C_0(X)$ 与 $\varepsilon>0$，定义
\[
 \mathcal N(\mu;G,\varepsilon)
 =\left\{\nu:
   \left|\int f\dd(\nu-\mu)\right|<\varepsilon
   \text{ 对每个 }f\in G
 \right\}.
\]
这些集合构成 $\mu$ 的弱星基本邻域。
\end{definition}

\begin{definition}[经验测度]\label{def:7-4-2-3}
若 $X_1,X_2,\ldots$ 为 $X$-值随机元，定义
\[
 L_n=\frac1n\sum_{k=1}^n\delta_{X_k}.
\]
它在每条样本路径上都是概率测度，并满足
$\int f\dd L_n=n^{-1}\sum_{k=1}^nf(X_k)$。
\end{definition}

\begin{theorem}[Dirac 组合弱星稠密]\label{theorem:7-4-2-1}
设 $X$ 为局部紧 Hausdorff 空间，$\mu$ 为有限符号或复 Radon 测度。存在有限原子测度网
$(\nu_\alpha)$，使
\[
 \nu_\alpha\weakstarto\mu,
 \qquad |\nu_\alpha|(X)\le|\mu|(X).
\]
网还可保持总质量 $\nu_\alpha(X)=\mu(X)$；若 $\mu\ge0$，可令每个 $\nu_\alpha\ge0$，从而同时保持
正性和总质量。若 $C_0(X)$ 可分，则可以用序列代替网。
\end{theorem}

\begin{proof}
令
\[
 \mathfrak D=\{(G,\varepsilon):G\subset C_0(X)\text{ 有限},\ \varepsilon>0\},
\]
并规定
$(G,\varepsilon)\preceq(H,\eta)$ 当且仅当 $G\subset H$ 且 $\eta\le\varepsilon$。它是有向集。

对每个 $\alpha=(G,\varepsilon)$，现在直接构造 $\nu_\alpha$。记 $M=|\mu|(X)$；若 $M=0$，取
$\nu_\alpha=0$。以下设 $M>0$，并选 $0<\delta<\varepsilon/(2M)$。由 $|\mu|$ 的 Radon 内正则性，
存在紧集 $K_0$ 使 $|\mu|(X\setminus K_0)<\delta$。由于 $G$ 有限且每个 $f\in G$ 在无穷远消失，
存在紧集 $K_1$ 使
\[
 |f(x)|<\delta\qquad(x\notin K_1,\ f\in G).
\]
令 $K=K_0\cup K_1$。对每个 $x\in K$，有限多个函数的连续性给出一个开邻域 $U_x$，使
\[
 |f(y)-f(z)|<\delta\qquad(y,z\in U_x,\ f\in G).
\]
从这些邻域取有限子覆盖 $U_1,\ldots,U_N$，并把 $K$ 分成 Borel 片
\[
 A_1=K\cap U_1,\qquad
 A_j=K\cap\left(U_j\setminus\bigcup_{i<j}U_i\right)\ (2\le j\le N),
 \qquad A_0=X\setminus K.
\]
删去空片，在每个余下的 $A_j$ 中选 $x_j$，并置
\[
 \nu_\alpha=\sum_j\mu(A_j)\delta_{x_j}.
\]
若 $A_0=\varnothing$，以下所有下标为 $0$ 的项均约定删去。
对 $f\in G$，有限分割给
\begin{align*}
 \left|\int f\,\dd(\nu_\alpha-\mu)\right|
 &\leq\sum_{j\geq1}\int_{A_j}|f(x_j)-f(x)|\,|\mu|(\dd x)
      +\int_{A_0}|f(x_0)-f(x)|\,|\mu|(\dd x)\\
 &\leq\delta|\mu|(K)+2\delta|\mu|(A_0)
 \leq2\delta M<\varepsilon,
\end{align*}
并且
\[
 |\nu_\alpha|(X)\leq\sum_j|\mu(A_j)|\leq|\mu|(X),
 \qquad
 \nu_\alpha(X)=\sum_j\mu(A_j)=\mu(X).
\]
若 $\mu\geq0$，所有系数都非负，所以 $\nu_\alpha\geq0$。

固定 $f\in C_0(X)$ 和 $\delta>0$ 后，所有
$\alpha\succeq(\{f\},\delta)$ 都满足
$|\int f\dd(\nu_\alpha-\mu)|<\delta$，这正是弱星收敛。

若 $C_0(X)$ 可分，取稠密列 $(f_j)$，令 $\nu_n$ 同时以误差 $1/n$ 匹配
$f_1,\ldots,f_n$。记 $M=|\mu|(X)$；当 $M=0$ 结论直接成立。给定 $f$，先选 $f_j$ 使
$\norm{f-f_j}_\infty<\delta/(4M)$。当 $n\ge j$ 时，
\begin{align*}
 \left|\int f\dd(\nu_n-\mu)\right|
 &\le\left|\int(f-f_j)\dd\nu_n\right|
   +\left|\int f_j\dd(\nu_n-\mu)\right|
   +\left|\int(f_j-f)\dd\mu\right|\\
 &\le2M\norm{f-f_j}_\infty+\frac1n<\delta
\end{align*}
对充分大 $n$ 成立。非可分情形没有一个可数稠密测试族可供这项对角化使用，故一般结论必须表述成网；
这并不排除某个特殊测度恰好拥有某条原子逼近序列。
\end{proof}

\begin{theorem}[有限测试的确定性构造]\label{theorem:7-4-2-2}
设 $\mu$ 为局部紧 Hausdorff 空间 $X$ 上的有限符号或复 Radon 测度。给定有限
$G\subset C_0(X)$ 与 $\varepsilon>0$，存在有限 Borel 分割
$X=A_0\dot\cup\cdots\dot\cup A_N$ 及点 $x_j\in A_j$（空片删去），使
\[
 \nu=\sum_{j=0}^N\mu(A_j)\delta_{x_j}
\]
满足
\begin{equation}\label{eq:7-4-2-finite-test-error}
 \left|\int f\dd(\nu-\mu)\right|<\varepsilon\quad(f\in G),
 \qquad |\nu|(X)\le|\mu|(X),
 \qquad \nu(X)=\mu(X).
\end{equation}
正测度情形中 $\nu$ 也为正测度。
\end{theorem}

\begin{proof}
令 $M=|\mu|(X)$。若 $M=0$，取 $\nu=0$。以下设 $M>0$，并取
$0<\eta<\varepsilon/(2M)$。Radon 内正则性给紧集 $K_0$，使
$|\mu|(X\setminus K_0)<\eta$。又因 $G$ 有限且每个函数在无穷远消失，存在紧集 $K_1$，使
\begin{equation}\label{eq:7-4-2-tail-small-functions}
 |f(x)|<\eta\qquad(x\notin K_1,\ f\in G).
\end{equation}
令 $K=K_0\cup K_1$。有限多个连续函数可在紧集 $K$ 上用同一个有限开覆盖控制振幅：对每个
$x\in K$ 可取开邻域
$U_x$，使
\begin{equation}\label{eq:7-4-2-common-oscillation}
 |f(y)-f(z)|<\eta
 \qquad(y,z\in U_x,\ f\in G).
\end{equation}
从 $(U_x)_{x\in K}$ 取有限子覆盖 $U_1,\ldots,U_N$，再令
\[
 A_1=K\cap U_1,\qquad
 A_j=K\cap\left(U_j\setminus\bigcup_{i<j}U_i\right),\qquad
 A_0=X\setminus K.
\]
删去空片，并在每个非空 $A_j$ 中选 $x_j$。这是有限 Borel 分割。若 $A_0\ne\varnothing$，
由 \eqref{eq:7-4-2-tail-small-functions}，$x_0$ 也使所有 $|f(x_0)|<\eta$。

对 $f\in G$，直接按分割计算
\begin{align*}
 \left|\int f\dd\nu-\int f\dd\mu\right|
 &=\left|\sum_{j=0}^N\int_{A_j}(f(x_j)-f(x))\dd\mu(x)\right|\\
 &\le\sum_{j=1}^N\eta|\mu|(A_j)+2\eta|\mu|(A_0)\\
 &\le2\eta M<\varepsilon.
\end{align*}
这里第一行只是有限和与积分的线性，没有交换无穷操作；第二行分别用了
\eqref{eq:7-4-2-common-oscillation} 和尾部两个函数值都小于 $\eta$。
此外，
\[
 |\nu|(X)\le\sum_{j=0}^N|\mu(A_j)|
 \le\sum_{j=0}^N|\mu|(A_j)=|\mu|(X),
\]
而 $\nu(X)=\sum_j\mu(A_j)=\mu(X)$。若 $\mu\ge0$，所有系数 $\mu(A_j)$ 非负，正性也被保留。
\end{proof}

在紧度量空间上，若 $G$ 中每个函数都是 $L$-Lipschitz，且分割片直径不超过 $h$，上述逐片计算立即给
\begin{equation}\label{eq:7-4-2-lipschitz-atomic-error}
 \left|\int f\dd(\nu-\mu)\right|
 \le Lh\,|\mu|(X).
\end{equation}
中点求积以网格直径控制误差，经验测度以大数律控制误差，而下一命题用有限维凸性把指定测试上的误差
直接降为零。

\begin{theorem}[经验测度强一致性]\label{theorem:7-4-2-3}
设 $X$ 为 Polish 空间，$X_1,X_2,\ldots$ 独立同分布，共同 Borel 概率分布为 $\mu$。则
\[
 L_n\Rightarrow\mu\qquad\text{a.s.}
\]
其中 $\Rightarrow$ 表示对所有有界连续函数的弱收敛。
\end{theorem}

\begin{proof}
先把所需的可数决定类真正构造出来。由 Polish 假设，取 $X$ 上一个兼容完备度量 $d$ 和一个可数
稠密集 $D=\{z_i:i\geq1\}$。对 $i\geq1$ 与有理数 $0<r<s$，定义帽函数
\[
 \psi_{i,r,s}(x)=
 \begin{cases}
  1,&d(x,z_i)\leq r,\\
  \dfrac{s-d(x,z_i)}{s-r},&r<d(x,z_i)<s,\\
  0,&d(x,z_i)\geq s.
 \end{cases}
\]
它取值于 $[0,1]$，且 Lipschitz 常数不超过 $(s-r)^{-1}$。令 $\mathcal F_0$ 由常数函数 $1$ 和任意
有限多个上述帽函数的逐点最大值组成；$\mathcal F_0$ 可数，并且其中每个函数都是有界 Lipschitz
函数。

这个固定族决定概率测度的弱收敛。事实上，设概率 $\rho_n$ 对 $\mathcal F_0$ 中每个函数的积分都收敛
到 $\rho$ 的相应积分。空开集的 Portmanteau 不等式没有内容。给定非空开集 $U$ 与 $x\in U$，若
$U\neq X$，记 $a=d(x,U^c)>0$，先取 $z_i$ 使 $d(x,z_i)<a/3$，再选有理数
$r<s$，使
\[
 d(x,z_i)<r<s<d(z_i,U^c).
\]
这样的有理数存在，因为
$d(z_i,U^c)\geq a-d(x,z_i)>2a/3>d(x,z_i)$。
于是 $\psi_{i,r,s}(x)=1$ 且 $\supp\psi_{i,r,s}\subset U$。因此，所有支撑包含于 $U$ 的
$\mathcal F_0$-函数可以枚举为 $(h_j)$；令 $H_m=h_1\vee\cdots\vee h_m$，则
$0\leq H_m\uparrow\1_U$。当 $U=X$ 时直接取 $H_m=1$。对每个固定 $m$，
\[
 \liminf_{n\to\infty}\rho_n(U)
 \geq\lim_{n\to\infty}\int H_m\dd\rho_n
 =\int H_m\dd\rho.
\]
令 $m\to\infty$ 并用单调收敛，得到
$\liminf_n\rho_n(U)\geq\rho(U)$。Portmanteau 定理~\ref{theorem:5-3-3-1} 因而给
$\rho_n\Rightarrow\rho$。这就由上述帽函数构造证明了所需的可数有界 Lipschitz 决定性质。

枚举 $\mathcal F_0=\{f_m:m\geq1\}$。固定 $m$。$f_m$ 有界，故 $f_m(X_k)$ i.i.d. 且可积。
强大数律给概率一事件 $\Omega_m$，使
\[
 \int f_m\dd L_n
 =\frac1n\sum_{k=1}^nf_m(X_k)
 \longrightarrow\E f_m(X_1)=\int f_m\dd\mu.
\]
可数交 $\Omega_0=\bigcap_{m\ge1}\Omega_m$ 仍有概率一。在每个
$\omega\in\Omega_0$ 上，上式对全部 $m$ 同时成立；收敛决定类的定义遂给
$L_n(\omega)\Rightarrow\mu$。

在 $X=\R$ 时，取半轴指标的连续上下夹逼，可把同一结论写成经验分布函数
$F_n(t)=n^{-1}\sum_k\1_{\{X_k\le t\}}$ 在 $F(t)=\mu(( -\infty,t])$ 的每个连续点收敛。
这里仍只需处理可数个决定点，再由单调性夹逼其余连续点；不能对不可数个 $t$ 直接取概率一事件的交。
\end{proof}

\begin{proposition}[Tchakaloff 有限矩匹配]\label{proposition:7-4-2-4}
设 $K$ 为紧 Hausdorff 空间，$\mu$ 为其上的有限正 Radon 测度，$V\subset C(K;\R)$ 为 $m$ 维线性空间。
则存在正原子测度
\[
 \nu=\sum_{j=1}^rc_j\delta_{x_j},\qquad
 0\leq r\leq m+1,
\]
使
\[
 \nu(K)=\mu(K),\qquad
 \int_Kf\dd\nu=\int_Kf\dd\mu\quad(f\in V).
\]
当 $r>0$ 时各 $c_j>0$；$r=0$ 表示空和，即零测度。
若 $1\in V$，可改进为 $r\le m$。这两个一般界都是最佳的。
\end{proposition}

\begin{proof}
若 $\mu(K)=0$，取 $r=0$ 的空和即得零测度。以下令 $M=\mu(K)>0$，并取 $V$ 的基
$f_1,\ldots,f_m$。把总质量坐标一并加入，定义连续映射
\[
 \Phi:K\to\R^{m+1},\qquad
 \Phi(x)=(f_1(x),\ldots,f_m(x),1),
\]
以及概率 $P=M^{-1}\mu$。向量
\[
 b=\int_K\Phi(x)\dd P(x)
\]
属于 $S:=\Phi(K)$ 的闭凸包。下面同时证明所需的凸包闭性，并给出有限维 \Caratheodory{} 约化步骤。
令 $\tau=\Phi_*P$。对每个 $n$，用有限个半径小于 $1/(2n)$ 的 Euclidean 球覆盖
紧集 $S$，再逐次作差得到有限 Borel 分割
$S=E_{n,1}\dot\cup\cdots\dot\cup E_{n,q_n}$，其中每个非空分片的直径小于 $1/n$。在每片选
$y_{n,j}\in E_{n,j}$，置
\[
 b_n=\sum_j\tau(E_{n,j})y_{n,j}\in\operatorname{conv}S.
\]
由于
\[
 \norm{b-b_n}
 =\left\|\sum_j\int_{E_{n,j}}(y-y_{n,j})\,\tau(\dd y)\right\|
 \leq\frac1n,
\]
故 $b\in\overline{\operatorname{conv}S}$。

再证明这个凸包已经闭。记 $d$ 为 $S$ 的仿射维数；由于 $S$ 落在末坐标等于 $1$ 的仿射超平面内，
$d\leq m$。若一个凸组合 $y=\sum_{j=1}^r\theta_jy_j$ 使用 $r>d+1$ 个正权重，则这些点仿射相关，
故存在不全为零的实数 $a_j$，使
\[
 \sum_{j=1}^ra_j=0,
 \qquad
 \sum_{j=1}^ra_jy_j=0.
\]
把 $(a_j)$ 乘以 $-1$（若有必要），可设某个 $a_j>0$。取
\[
 t=\min_{a_j>0}\frac{\theta_j}{a_j}>0,
 \qquad \theta_j'=\theta_j-ta_j.
\]
则每个 $\theta_j'\geq0$，至少一个为零，并且
$\sum_j\theta_j'=1$、$\sum_j\theta_j'y_j=y$。反复消去零权，便知每个凸包中的点都能用至多
$d+1$ 个点表示。于是 $\operatorname{conv}S$ 是紧集
\[
 S^{d+1}\times
 \left\{(\theta_1,\ldots,\theta_{d+1})\in[0,1]^{d+1}:\sum_j\theta_j=1\right\}
\]
在连续映射 $(y_1,\ldots,y_{d+1},\theta)\mapsto\sum_j\theta_jy_j$ 下的像，因而为闭集。
所以 $b\in\operatorname{conv}S$，并存在 $x_1,\ldots,x_r\in K$ 和 $\theta_j\geq0$，使
\[
 r\leq d+1\leq m+1,\qquad \sum_j\theta_j=1,\qquad
 b=\sum_{j=1}^r\theta_j\Phi(x_j).
\]
删去零权并令 $c_j=M\theta_j$，便同时得到总质量和基函数积分的匹配；线性性把它推广到全部 $V$。

若 $1\in V$，可选 $f_1=1$。这在“末坐标等于 $1$”之外再给出关系 $y_1=y_{m+1}$，故
$\Phi(K)$ 的仿射维数至多 $m-1$。把上面的仿射相关消元用于这个较小的仿射空间，只需至多 $m$ 个点。

最后核对界的准确性。不假定 $1\in V$ 时，令 $K=\{0,e_1,\ldots,e_m\}\subset\R^m$，令 $V$ 由
$m$ 个坐标函数在 $K$ 上的限制张成，并让 $\mu$ 在这 $m+1$ 个点上都给正质量。归一化积分向量位于
单纯形内部，不能由少于 $m+1$ 个顶点的凸组合表示。若 $1\in V$，取含 $m$ 个点的有限空间
$K$、$V=C(K)$ 及在每点都给正质量的 $\mu$；精确匹配所有 $V$ 中函数会逐点匹配质量，故必须使用
全部 $m$ 个原子。
\end{proof}

\begin{figure}[htbp]
\centering
\begin{tikzpicture}[node distance=11mm and 15mm]
  \node[book strong node] (space) {空间 $K$\\点 $x$};
  \node[book strong node,right=of space] (image) {有限观测像\\$\Phi(x)\in\R^m$};
  \node[book warm node,right=of image] (bar) {测度重心\\$b=\int\Phi\dd P$};
  \node[book strong node,below=of bar] (atom) {有限凸组合\\$\sum_{j=1}^r\theta_j\Phi(x_j)$};
  \node[book warm node,left=of atom] (emp) {经验重心\\$n^{-1}\sum_{k=1}^n\Phi(X_k)$};
  \draw[book accent arrow] (space)--node[above,book label]{$m$ 个测试}(image);
  \draw[book accent arrow] (image)--node[above,book label]{积分}(bar);
  \draw[book accent arrow] (atom)--node[right,book label]{精确匹配}(bar);
  \draw[book dashed arrow] (emp)--node[below,book label]{大数律}(atom);
  \draw[book dashed arrow] (space)--node[left,book label]{抽样}(emp);
\end{tikzpicture}
\caption{有限测试把测度问题降到有限维重心问题。确定性分割近似重心，经验测度随机逼近重心，Tchakaloff 用有限凸组合精确表示重心。}
\label{fig:7-4-2-finite-observation-barycenter}
\end{figure}

图~\ref{fig:7-4-2-finite-observation-barycenter} 只在有限测试像中宣称精确。若不断扩大测试族，必须另用
网、可数决定类或稠密性完成极限；一次 Tchakaloff 表示绝不意味着原测度本身已经是有限原子的。

\subsection{结果汇总：由实分析通往对偶空间、矩问题与谱理论}\label{subsec:7-4-3}
\index{结果汇总}\index{Feller 半群}\index{生成元}\index{预解算子}

前七章反复出现同一个四阶段结构：有限测试先给一致的数据，完备性、扩张或紧性制造无限对象，稠密类
或变换负责识别，核与半群再描述对象如何演化。下面先用三个例子说明从测试信息读取
表示对象时不可省略的假设，再由三道综合题分别展开 Fourier 完备化、矩泛函表示和随机半群这
三条路线。一般弱星紧性、闭算子与谱测度会在后续泛函分析中推广这些构造。

\begin{example}[变分到测度]\label{ex:7-4-3-1}
设 $a<x_0<b$，$f,h\in C^1[a,b]$，$\Phi\in C^1(\R)$。设
$L:[a,b]\times\R^2\to\R$ 连续，并且对后两个变量的偏导数 $L_u,L_p$ 存在且连续。对
\[
 S(f)=\Phi(f(x_0))+\int_a^bL(x,f(x),f'(x))\dd x
\]
考虑 $f+\varepsilon h$。映射
\[
 (x,\varepsilon)\longmapsto
 \bigl(x,f(x)+\varepsilon h(x),f'(x)+\varepsilon h'(x)\bigr)
\]
把紧集 $[a,b]\times[-\varepsilon_0,\varepsilon_0]$ 送到紧集，因而 $L_u,L_p$ 在该像上有共同有限上界；
差商由 $C(|h|+|h'|)$ 支配。支配收敛定理给
\begin{equation}\label{eq:7-4-3-first-variation}
 \delta S(f)[h]
 =\Phi'(f(x_0))h(x_0)
 +\int_a^b\bigl(L_u(x,f,f')h+L_p(x,f,f')h'\bigr)\dd x.
\end{equation}
若 $x\mapsto L_p(x,f(x),f'(x))$ 为 $C^1$，分部积分得到
\begin{align}\label{eq:7-4-3-variation-measure}
 \delta S(f)[h]
 &=\int_a^b\left(L_u-\frac{\dd}{\dd x}L_p\right)h\dd x
 +L_p(b,f(b),f'(b))h(b)-L_p(a,f(a),f'(a))h(a)\notag\\
 &\quad+\Phi'(f(x_0))h(x_0).
\end{align}
因此该一阶变分在上确界范数下由下列符号测度表示：
\[
 \nu=\left(L_u-\frac{\dd}{\dd x}L_p\right)\dd x
      -L_p(a,f(a),f'(a))\delta_a
      +L_p(b,f(b),f'(b))\delta_b
      +\Phi'(f(x_0))\delta_{x_0},
\]
其中密度中的 $L_u,L_p$ 均沿曲线
$x\mapsto(x,f(x),f'(x))$ 取值；式~\eqref{eq:7-4-3-variation-measure} 正是
$\delta S(f)[h]=\int h\dd\nu$。若只知道 \eqref{eq:7-4-3-first-variation} 而不能对 $L_p$ 分部积分，
含 $h'$ 的项只在 $C^1$ 范数下连续，未必由测度表示。例~\ref{ex:7-4-1-3} 的导数泛函正是这个边界。
\end{example}

\begin{example}[矩到弱星极限]\label{ex:7-4-3-2}
设 $K=[-1,1]$，概率测度 $\mu_N,\mu$ 都支持在 $K$，并且二者的 $0$ 到 $N$ 阶矩相同。给定
$f\in C(K)$ 和
$\varepsilon>0$，Stone--Weierstrass 定理给多项式 $p$，使
$\norm{f-p}_\infty<\varepsilon$；若 $N\ge\deg p$，则
\[
 \left|\int f\dd\mu_N-\int f\dd\mu\right|
 \le\int|f-p|\dd\mu_N+\int|f-p|\dd\mu<2\varepsilon.
\]
所以共同紧支撑和总质量控制把逐阶矩匹配升级为弱收敛。若删去共同紧支撑，仅有所有 Hankel 二次型非负
不能给紧性；例如 $\delta_N$ 的质量直接逃向无穷。
\end{example}

\begin{example}[半群到生成元]\label{ex:7-4-3-3}
在 $\R^d$ 上令 $P_tf(x)=\E f(x+B_t)$，其中 $B_t\sim N(0,tI_d)$。对 Fourier 模式
$e_\xi(x)=e^{i\xi\cdot x}$，Gaussian 特征函数给
\[
 P_te_\xi(x)=e^{-t|\xi|^2/2}e_\xi(x),
 \qquad
 \frac{P_te_\xi-e_\xi}{t}\longrightarrow-\frac{|\xi|^2}{2}e_\xi
 =\frac12\Delta e_\xi.
\]
这项计算确定了 Fourier 符号 $-|\xi|^2/2$，但 $e_\xi\notin C_0(\R^d)$，所以它不能替代生成元定义域的证明。
要把符号识别为生成元，还必须在生成元的定义域上证明上确界范数极限。综合题
\ref{problem:7-4-3-4} 对 $C_c^\infty(\R^d)$ 中的函数完成这一步。
\end{example}

\begin{remark}[假设的必要性]
表~\ref{tab:7-4-3-assumptions-counterexamples} 汇总了七个基本反例或失效情形。它们分别说明单调连续性、
局部紧性、连续性、可分性、可数决定类、支集正性和 Feller 条件为何出现在相应结论中。
\end{remark}

\begin{table}[htbp]
\centering
\small
\begin{tabularx}{\textwidth}{>{\raggedright\arraybackslash}p{0.19\textwidth}
 >{\raggedright\arraybackslash}p{0.22\textwidth}
 >{\raggedright\arraybackslash}p{0.31\textwidth}
 >{\raggedright\arraybackslash}X}
\toprule
结论 & 关键假设 & 反例或失效情形 & 相关主题 \\
\midrule
Daniell--Stone & $f_n\downarrow0\Rightarrow I(f_n)\downarrow0$
& 把收敛数列空间上的通常极限用 Hahn--Banach 保范延拓为
$I:\ell^\infty(\N)\to\R$；此延拓满足 $I(\mathbf 1)=1$。若 $0\le x\le M\mathbf 1$，则
$\|M\mathbf 1-x\|_\infty\le M$，从而
$I(M\mathbf 1-x)\le |I(M\mathbf 1-x)|\le M$，并有
$I(x)=M-I(M\mathbf 1-x)\ge0$，故 $I$ 正；但
$I(\1_{\{k\ge n\}})=1$，而这些函数逐点下降到零，故相应有限加法不能升级为可数加法
& 测度扩张、概率分布 \\
Riesz--Markov & 局部紧 Hausdorff 与局部 Urysohn 帽函数
& 没有紧开夹逼时，开集值的局部有限性与内正则性均无由推出
& $C_0(X)^*$、Radon 对偶 \\
测度表示导数 & $C_0$ 上确界范数连续
& $f_n(x)=n^{-1}\phi(nx)$ 使 $\norm{f_n}_\infty\to0$ 而 $f_n'(0)$ 不变
& 分布与 Sobolev 对偶 \\
原子逼近序列 & $C_0(X)$ 可分
& 非可分时弱星邻域没有可数稠密测试族；一般密度结论只能由有向网表达
& 弱星拓扑 \\
经验测度弱收敛 & Polish 空间的可数决定类
& 对不可数测试族分别取概率一事件，所得不可数交未必仍有概率一
& Varadarajan、经验过程 \\
$K$-支撑矩表示 & 对所有在 $K$ 上非负的多项式为正
& $K=[-1,1]$ 时 $L(p)=p(2)$ 满足 $L(q^2)\ge0$，却有
$L(1-x^2)=-3<0$；Hankel 正性不能保证表示测度支撑于 $K$
& Hausdorff 矩问题 \\
Feller 半群 & $K(C_0)\subset C_0$
& $Q(x)=\delta_0$（$x\in\Q$）、$Q(x)=\delta_1$（$x\notin\Q$）使
$Kf$ 对 $f(0)\ne f(1)$ 不连续
& 生成元理论 \\
\bottomrule
\end{tabularx}
\caption{结论、关键假设与反例。第三列说明每项假设在相应构造中所承担的作用。}
\label{tab:7-4-3-assumptions-counterexamples}
\end{table}

\begin{figure}[htbp]
\centering
\begin{tikzpicture}[node distance=8mm and 17mm]
  \node[book strong node] (obs) {有限观测\\邻域、测试、矩、有限维分布};
  \node[book strong node,below=of obs] (exist) {存在\\完备化、扩张、紧性};
  \node[book strong node,below=of exist] (id) {识别\\稠密、微分、变换};
  \node[book strong node,below=of id] (evo) {演化\\概率核、半群、预解算子};
  \node[book warm node,right=of exist] (dual) {对偶空间\\$C_0^*$、矩泛函};
  \node[book warm node,right=of evo] (spectral) {算子方向\\生成元、预解式};
  \draw[book accent arrow] (obs)--node[right,book label]{一致有限数据}(exist);
  \draw[book accent arrow] (exist)--node[right,book label]{候选对象}(id);
  \draw[book accent arrow] (id)--node[right,book label]{唯一结构}(evo);
  \draw[book dashed arrow] (obs.east)--(dual.north);
  \draw[book dashed arrow] (id.east)--(dual.south);
  \draw[book dashed arrow] (evo)--(spectral);
\end{tikzpicture}
\caption{有限观测到表示与演化的四层结构。三项综合题分别展开 Fourier 完备化、矩泛函表示和随机半群三条路线。}
\label{fig:7-4-3-interface-master}
\end{figure}

图~\ref{fig:7-4-3-interface-master} 中“存在”和“识别”不能互换：先有一族相容 Fourier 部分和，仍须
$L^2$ 完备性产生极限；先有矩截断解，仍须紧支撑和稠密多项式识别极限；先有转移核，仍须
$C_0$ 不变性与强连续性，才谈得上该 Banach 空间中的生成元。

\begin{problem}[综合题 A：三角矩与 $L^2$ 完备化]\label{problem:7-4-3-2}
给定 $(a_k)_{k\in\Z}\in\ell^2(\Z)$。仅用有限 Fourier 和、正交性、\Fejer{} 逼近与 $L^2$ 完备性，
证明存在唯一 $f\in L^2(\T)$ 使 $\widehat f(k)=a_k$ 对所有 $k\in\Z$ 成立，并证明
$\norm f_2^2=\sum_k|a_k|^2$。
\end{problem}

\begin{proof}[解答]
在 $\T=\R/(2\pi\Z)$ 上使用归一化 Haar 测度 $\dd m=\dd x/(2\pi)$，并置
$e_k(x)=e^{ikx}$。直接积分给
\[
 \langle e_j,e_k\rangle
 =\frac1{2\pi}\int_{-\pi}^{\pi}e^{i(j-k)x}\dd x
 =\begin{cases}1,&j=k,\\0,&j\ne k.\end{cases}
\]
定义有限和
\[
 S_N=\sum_{|k|\le N}a_ke_k.
\]
若 $M>N$，有限正交展开给出精确恒等式
\begin{equation}\label{eq:7-4-3-fourier-cauchy}
 \norm{S_M-S_N}_2^2=\sum_{N<|k|\le M}|a_k|^2.
\end{equation}
因为 $(a_k)\in\ell^2$，右端随 $M,N\to\infty$ 趋于零。因此 $(S_N)$ 是 $L^2(\T)$ 中的 Cauchy
列。$L^2$ 完备性给出 $f\in L^2(\T)$，使 $S_N\to f$ 于 $L^2$。这正是由有限 Fourier 和产生
极限函数的存在性步骤。

固定 $j\in\Z$。Cauchy--Schwarz 不等式和 $\norm{e_j}_2=1$ 给
\[
 |\widehat f(j)-\widehat{S_N}(j)|
 =|\langle f-S_N,e_j\rangle|
 \le\norm{f-S_N}_2\longrightarrow0.
\]
当 $N\ge|j|$ 时 $\widehat{S_N}(j)=a_j$，所以 $\widehat f(j)=a_j$。再由范数连续性与有限
Parseval 恒等式，
\[
 \norm f_2^2=\lim_{N\to\infty}\norm{S_N}_2^2
 =\lim_{N\to\infty}\sum_{|k|\le N}|a_k|^2
 =\sum_{k\in\Z}|a_k|^2.
\]

只剩唯一性。若 $g\in L^2(\T)$ 的全部 Fourier 系数为零，则它的 \Fejer{} 均值为
\[
 \sigma_Ng
 =\sum_{|k|\le N}\left(1-\frac{|k|}{N+1}\right)\widehat g(k)e_k=0.
\]
定理~\ref{theorem:7-2-1-1} 的 \Fejer{} 近似先对连续函数给一致收敛；连续函数在 $L^2$ 中稠密，
而卷积正核满足 $\norm{\sigma_Nh}_2\le\norm h_2$。这里的收缩并非另行调用 Plancherel：若
$K_N\geq0$ 是归一化 \Fejer{} 核，则 Jensen 不等式逐点给
\[
 |\sigma_Nh(x)|^2
 =\left|\int_\T K_N(y)h(x-y)\dd m(y)\right|^2
 \leq\int_\T K_N(y)|h(x-y)|^2\dd m(y).
\]
对 $x$ 积分并用 Tonelli 以及 Haar 测度的平移不变性，得到
$\norm{\sigma_Nh}_2^2\leq\norm h_2^2$。所以一个三项逼近给
$\sigma_Ng\to g$ 于 $L^2$。具体地，给定连续 $q$，
\[
 \norm{\sigma_Ng-g}_2
 \le\norm{\sigma_N(g-q)}_2+\norm{\sigma_Nq-q}_2+\norm{q-g}_2
 \le2\norm{g-q}_2+\norm{\sigma_Nq-q}_\infty.
\]
先选 $q$，再令 $N\to\infty$，右端趋于零。因此 $g=0$。若两个函数具有同一系数，对其差应用此论证
便得二者在 $L^2$ 中相同。

综上，正交性给出 Cauchy 估计~\eqref{eq:7-4-3-fourier-cauchy}，$L^2$ 完备性产生极限，
系数泛函的连续性恢复给定数据，\Fejer{} 逼近则给出唯一性。
\end{proof}

\begin{problem}[综合题 B：多项式矩与 Stieltjes 变换]\label{problem:7-4-3-3}
固定紧区间 $K=[a,b]$。设 $L:\R[x]\to\R$ 线性，并且每个在 $K$ 上非负的实多项式 $p$ 都满足
$L(p)\ge0$。证明 $L$ 由 $K$ 上唯一有限正测度 $\mu$ 表示；构造在 $K$ 上匹配 $0$ 到 $N$ 阶矩的
有限原子测度 $\mu_N$；证明 $\mu_N\Rightarrow\mu$，并证明其 Stieltjes 变换在
$\C\setminus K$ 上局部一致收敛。
\end{problem}

\begin{proof}[解答]
\textbf{第一步：把多项式泛函延到 $C(K)$。}
正性给 $L(1)\ge0$。对任意实多项式 $p$，令 $M_p=\norm p_{C(K)}$。多项式
$M_p\pm p$ 在 $K$ 上非负，故
\[
 0\le L(M_p\pm p)=M_pL(1)\pm L(p).
\]
于是
\begin{equation}\label{eq:7-4-3-polynomial-bound}
 |L(p)|\le L(1)\norm p_{C(K)}.
\end{equation}
若 $L(1)=0$，该式给 $L=0$，以下结论取零测度即可。设 $L(1)>0$。Stone--Weierstrass 定理说明
实多项式在 $C(K;\R)$ 中一致稠密；估计 \eqref{eq:7-4-3-polynomial-bound} 因而给唯一连续延拓
$\widetilde L:C(K)\to\R$，且 $\norm{\widetilde L}\le L(1)$。

还须验证延拓保持正性。给 $f\in C(K)$、$f\ge0$，取多项式 $p_n\to f$ 一致，并令
$\varepsilon_n=\norm{p_n-f}_\infty$。则 $p_n+\varepsilon_n\ge0$ 于 $K$，所以
\[
 L(p_n)+\varepsilon_nL(1)\ge0.
\]
令 $n\to\infty$ 得 $\widetilde L(f)\ge0$。定理~\ref{theorem:7-4-1-2} 应用于紧空间 $K$，得到唯一
有限正 Radon 测度 $\mu$，满足
\[
 L(p)=\int_Kp(x)\dd\mu(x),\qquad \mu(K)=\widetilde L(1)=L(1).
\]

\textbf{第二步：Tchakaloff 原子截断。}
对 $N\ge0$ 令
$V_N=\Span\{1,x,\ldots,x^N\}$。它的维数为 $N+1$ 且含常数。命题
\ref{proposition:7-4-2-4} 给出支撑仍在 $K$ 的正原子测度
\[
 \mu_N=\sum_{j=1}^{r_N}c_{N,j}\delta_{x_{N,j}},
 \qquad r_N\le N+1,\quad c_{N,j}>0,\quad x_{N,j}\in K,
\]
使
\begin{equation}\label{eq:7-4-3-moment-match}
 \mu_N(K)=\mu(K),\qquad
 \int_Kx^k\dd\mu_N=\int_Kx^k\dd\mu=L(x^k)\quad(0\le k\le N).
\end{equation}
这里的 $N+1$ 而不是 $N+2$，来自 $1\in V_N$ 所降低的一维仿射维数。

\textbf{第三步：弱收敛。}
记 $M=L(1)$。给定 $f\in C(K)$ 和 $\varepsilon>0$，选多项式 $p$ 使
$\norm{f-p}_\infty<\varepsilon/(2M+1)$。当 $N\ge\deg p$ 时，
\eqref{eq:7-4-3-moment-match} 和线性性给 $\int p\dd\mu_N=\int p\dd\mu$，从而
\begin{align*}
 \left|\int f\dd\mu_N-\int f\dd\mu\right|
 &\le\int|f-p|\dd\mu_N+\int|f-p|\dd\mu\\
 &\le2M\norm{f-p}_\infty<\varepsilon.
\end{align*}
故 $\mu_N\Rightarrow\mu$。共同紧支集使紧性自动成立；证明甚至不需要抽取子列。

\textbf{第四步：Stieltjes 变换局部一致收敛。}
采用
\[
 S_\rho(z)=\int_K\frac1{z-x}\dd\rho(x),\qquad z\in\C\setminus K
\]
的符号。固定紧集 $C\Subset\C\setminus K$，令
$d=\dist(C,K)>0$。核满足共同界
\begin{equation}\label{eq:7-4-3-stieltjes-kernel-bounds}
 \left|\frac1{z-x}\right|\le d^{-1},\qquad
 \left|\frac1{z-x}-\frac1{w-x}\right|
 \le d^{-2}|z-w|\quad(z,w\in C,x\in K).
\end{equation}
第一项是所有积分存在以及参数积分的支配；第二项给
\[
 |S_{\mu_N}(z)-S_{\mu_N}(w)|\le Md^{-2}|z-w|,
 \qquad
 |S_\mu(z)-S_\mu(w)|\le Md^{-2}|z-w|.
\]
对每个固定 $z\in C$，函数 $x\mapsto(z-x)^{-1}$ 连续，弱收敛给
$S_{\mu_N}(z)\to S_\mu(z)$。给定 $\varepsilon>0$，取 $C$ 的有限
$\eta$-网 $z_1,\ldots,z_q$，其中 $2Md^{-2}\eta<\varepsilon/2$。在这有限多个点上同时取足够大的
$N$，使差小于 $\varepsilon/2$。对任意 $z\in C$ 选 $z_j$ 满足 $|z-z_j|<\eta$，于是
\begin{align*}
 |S_{\mu_N}(z)-S_\mu(z)|
 &\le|S_{\mu_N}(z)-S_{\mu_N}(z_j)|
 +|S_{\mu_N}(z_j)-S_\mu(z_j)|\\
 &\quad+|S_\mu(z_j)-S_\mu(z)|<\varepsilon.
\end{align*}
所以收敛在每个 $C\Subset\C\setminus K$ 上一致。

最后说明正性假设为何必须精确写成“在 $K$ 上非负”。仅有 Hankel 正性
$L(q^2)\ge0$ 不足以保证表示测度支撑于 $K$：在 $K=[-1,1]$ 上取 $L(p)=p(2)$，所有平方都取非负值，但
$1-x^2\ge0$ 于 $K$ 而 $L(1-x^2)=-3$。相应表示测度是 $\delta_2$，支撑恰在 $K$ 外。
\end{proof}

\begin{problem}[综合题 C：随机演化与算子方法]\label{problem:7-4-3-4}
设 $X$ 为局部紧 Hausdorff 空间，$(Q_t)_{t\ge0}$ 是满足 Chapman--Kolmogorov 方程的 Borel 转移核族，
$Q_0(x,\cdot)=\delta_x$，并令
$P_tf(x)=\int f(y)Q_t(x,\dd y)$。假定 $P_tC_0(X)\subset C_0(X)$ 且
$\norm{P_tf-f}_\infty\to0$（$t\downarrow0$）对每个 $f\in C_0(X)$ 成立。证明 $(P_t)$ 是
$C_0(X)$ 上的正收缩强连续半群；计算 Brownian 半群在 $C_c^\infty(\R^d)$ 上的生成元，并计算
满足 Feller 条件的 Poisson 跳半群在整个 $C_0(X)$ 上的生成元；最后对任意 $\alpha,\beta>0$ 证明
Laplace 预解算子恒等式
\[
 R_\alpha-R_\beta=(\beta-\alpha)R_\alpha R_\beta,
 \qquad R_\alpha f=\int_0^\infty e^{-\alpha t}P_tf\dd t.
\]
\end{problem}

\begin{proof}[解答]
\textbf{第一步：从转移族得到 Feller 半群。}
对 $f\in C_0(X)$，概率核给
\[
 |P_tf(x)|\le\int|f(y)|Q_t(x,\dd y)\le\norm f_\infty,
\]
故 $P_t$ 为正收缩。Chapman--Kolmogorov 及有界函数的核积分结合律给
\begin{align*}
 P_sP_tf(x)
 &=\int_X\int_Xf(z)Q_t(y,\dd z)Q_s(x,\dd y)\\
 &=\int_Xf(z)Q_{s+t}(x,\dd z)=P_{s+t}f(x).
\end{align*}
两重积分的绝对值由 $\norm f_\infty$ 支配，概率核上的 Tonelli/Fubini 定理保证交换合法。
$P_0=I$。零时刻的强连续性已假定；收缩性把它传播到每个时刻。对 $h\downarrow0$，
\[
 \norm{P_{t+h}f-P_tf}_\infty
 \le\norm{P_hf-f}_\infty\to0,
\]
而 $0<h<t$ 时
\[
 \norm{P_tf-P_{t-h}f}_\infty
 \le\norm{P_hf-f}_\infty\to0.
\]
因此 $(P_t)$ 是 $C_0(X)$ 上的正收缩强连续半群。其生成元定义为
\[
 D(A)=\left\{f\in C_0(X):
 \lim_{t\downarrow0}\frac{P_tf-f}{t}\text{ 在 }\norm{\cdot}_\infty\text{ 中存在}\right\},
 \qquad Af=\lim_{t\downarrow0}\frac{P_tf-f}{t}.
\]

\textbf{第二步：Brownian 生成元。}
在 $X=\R^d$ 上令 $P_tf(x)=\E f(x+\sqrt t Z)$，其中 $Z\sim N(0,I_d)$。取
$f\in C_c^\infty(\R^d)$。二阶 Taylor 公式的积分余项给
\begin{align*}
 f(x+y)-f(x)
 &=\nabla f(x)\cdot y+\frac12y^TD^2f(x)y+r(x,y),\\
 |r(x,y)|
 &\le\frac12|y|^2\omega_{D^2f}(|y|),
\end{align*}
其中
$\omega_{D^2f}(r)=\sup_{|x-x'|\le r}\norm{D^2f(x)-D^2f(x')}$。因为
$D^2f\in C_c(\R^d)$，它在整个 $\R^d$ 上一致连续，故该模量随 $r\downarrow0$ 趋于零。
Gaussian 一、二阶矩为
$\E Z_i=0$、$\E Z_iZ_j=\delta_{ij}$，所以
\[
 \frac{P_tf(x)-f(x)}t-\frac12\Delta f(x)
 =\frac1t\E r(x,\sqrt t Z).
\]
并且一致于 $x$ 有
\begin{equation}\label{eq:7-4-3-brownian-generator-domination}
 \frac1t|r(x,\sqrt t Z)|
 \le\frac12|Z|^2\omega_{D^2f}(\sqrt t|Z|)
 \le\norm{D^2f}_\infty|Z|^2.
\end{equation}
右端可积，左端对每个固定 $Z$ 趋于零。支配收敛定理应用于
\eqref{eq:7-4-3-brownian-generator-domination}，并在取期望前先取 $x$ 的上确界，得到
\[
 \left\|\frac{P_tf-f}{t}-\frac12\Delta f\right\|_\infty\longrightarrow0.
\]
故 $C_c^\infty(\R^d)\subset D(A)$ 且 $Af=\frac12\Delta f$。这一步给的是定义域内的范数极限，
不是只在 Fourier 模式上的形式符号。

\textbf{第三步：Poisson 跳生成元。}
设 $Q(x,\dd y)$ 为单次跳跃核，
$Kf(x)=\int f(y)Q(x,\dd y)$，并明确假设 Feller 条件
$K(C_0(X))\subset C_0(X)$。$K$ 是正收缩。若跳跃时钟速率为 $\lambda>0$，则
\begin{equation}\label{eq:7-4-3-poisson-semigroup-series}
 P_tf=e^{-\lambda t}\sum_{n=0}^\infty\frac{(\lambda t)^n}{n!}K^nf.
\end{equation}
该级数在 $C_0$ 范数中绝对收敛，因为
$\norm{K^nf}_\infty\le\norm f_\infty$。它确实定义强连续半群。对 $s,t\geq0$，双重级数满足
\[
 \sum_{m,n\geq0}e^{-\lambda(s+t)}
 \frac{(\lambda s)^m(\lambda t)^n}{m!n!}\norm{K^{m+n}f}_\infty
 \leq\norm f_\infty,
\]
故可在 $C_0$ 中按绝对收敛重排。令 $k=m+n$ 并用二项式公式，得到
\begin{align*}
 P_sP_tf
 &=e^{-\lambda(s+t)}
   \sum_{k=0}^\infty
   \left(\sum_{m=0}^k\frac{(\lambda s)^m(\lambda t)^{k-m}}{m!(k-m)!}\right)K^kf\\
 &=e^{-\lambda(s+t)}\sum_{k=0}^\infty
   \frac{(\lambda(s+t))^k}{k!}K^kf
 =P_{s+t}f.
\end{align*}
各项系数非负且和为一，所以 $P_t$ 是正收缩；并且
\[
 \norm{P_tf-f}_\infty
 \leq2(1-e^{-\lambda t})\norm f_\infty\longrightarrow0,
\]
故它在零时刻强连续。现在分离 $n=0,1$ 项得
\begin{align*}
 \frac{P_tf-f}{t}
 &=\frac{e^{-\lambda t}-1}{t}f
 +e^{-\lambda t}\lambda Kf
 +\frac{e^{-\lambda t}}t\sum_{n=2}^\infty\frac{(\lambda t)^n}{n!}K^nf.
\end{align*}
余项满足
\[
 \left\|\frac{e^{-\lambda t}}t
 \sum_{n=2}^\infty\frac{(\lambda t)^n}{n!}K^nf\right\|_\infty
 \le\frac{1-e^{-\lambda t}(1+\lambda t)}t\norm f_\infty\longrightarrow0.
\]
前两项分别趋于 $-\lambda f$ 与 $\lambda Kf$，故在整个 $C_0(X)$ 上
\[
 Af=\lambda(K-I)f.
\]
若只有 $Q$ 的可测性而没有 $K(C_0)\subset C_0$，式
\eqref{eq:7-4-3-poisson-semigroup-series} 仍可在有界可测函数上逐点有意义，却不能据此宣称得到
$C_0$ 上的半群；表~\ref{tab:7-4-3-assumptions-counterexamples} 已给出相应反例。

\textbf{第四步：Laplace 预解算子及预解式。}
设现在 $(P_t)$ 为第一步的一般正收缩强连续半群。对 $\alpha>0$，映射
$t\mapsto e^{-\alpha t}P_tf$ 在 $C_0(X)$ 中连续，且
\[
 \norm{e^{-\alpha t}P_tf}_\infty\le e^{-\alpha t}\norm f_\infty,
 \qquad \int_0^\infty e^{-\alpha t}\norm f_\infty\dd t
 =\frac{\norm f_\infty}{\alpha}.
\]
所以 Bochner 积分 $R_\alpha f$ 存在，并有
$\norm{R_\alpha}\le1/\alpha$。

对 $\alpha,\beta>0$，收缩估计给
\[
 \int_0^\infty\int_0^\infty
 e^{-\alpha s-\beta t}\norm{P_{s+t}f}_\infty\dd t\dd s
 \le\frac{\norm f_\infty}{\alpha\beta}<\infty.
\]
因此 Bochner--Fubini 定理允许交换积分；半群律给
\begin{align*}
 R_\alpha R_\beta f
 &=\int_0^\infty\int_0^\infty
 e^{-\alpha s-\beta t}P_{s+t}f\dd t\dd s\\
 &=\int_0^\infty P_uf
 \left(\int_0^ue^{-\alpha s-\beta(u-s)}\dd s\right)\dd u.
\end{align*}
若 $\alpha\ne\beta$，括号内等于
$(e^{-\alpha u}-e^{-\beta u})/(\beta-\alpha)$，故
\[
 R_\alpha R_\beta f=\frac{R_\alpha f-R_\beta f}{\beta-\alpha},
\]
即所需预解式；$\alpha=\beta$ 时原式两端都为零。

还可直接核对它确实解生成元方程。由有界算子与 Bochner 积分可交换以及半群律，
\begin{align*}
 P_hR_\alpha f
 &=\int_0^\infty e^{-\alpha t}P_{t+h}f\dd t\\
 &=e^{\alpha h}\left(R_\alpha f-
   \int_0^he^{-\alpha u}P_uf\dd u\right).
\end{align*}
于是
\[
 \frac{P_hR_\alpha f-R_\alpha f}{h}
 =\frac{e^{\alpha h}-1}{h}R_\alpha f
 -e^{\alpha h}\frac1h\int_0^he^{-\alpha u}P_uf\dd u.
\]
第一项在范数中趋于 $\alpha R_\alpha f$。第二项的平均在范数中趋于 $f$，因为
\[
 \left\|\frac1h\int_0^he^{-\alpha u}P_uf\dd u-f\right\|_\infty
 \le\frac1h\int_0^h\norm{e^{-\alpha u}P_uf-f}_\infty\dd u\longrightarrow0;
\]
被积函数由 $2\norm f_\infty$ 支配并由强连续性趋于零。因此
\[
 R_\alpha f\in D(A),\qquad
 AR_\alpha f=\alpha R_\alpha f-f,
\]
即 $(\alpha I-A)R_\alpha=I$。这一恒等式是 Hille--Yosida 理论中生成元与预解算子关系的初等形式。
\end{proof}

\begin{exercise}
在定理~\ref{theorem:7-4-1-2} 的证明中，逐项标出局部 Urysohn 引理、有限单位分解、
\Caratheodory{} 定理和层蛋糕公式分别负责哪一个箭头，并说明删去任一项后证明停在哪里。
\end{exercise}

\begin{exercise}
设 $X$ 为紧度量空间，$\mu$ 为概率测度，$f_1,\ldots,f_m$ 都是 $L$-Lipschitz。构造直径至多 $h$
的有限 Borel 分割，并从头证明 \eqref{eq:7-4-2-lipschitz-atomic-error}；比较所得原子数与
Tchakaloff 的 $m+1$ 上界为何控制的是不同量。
\end{exercise}

\begin{exercise}
对 Poisson 跳半群证明 $P_sP_t=P_{s+t}$：在
\eqref{eq:7-4-3-poisson-semigroup-series} 中逐项相乘，写出绝对收敛支配，再用二项式恒等式合并
$K^{m+n}$ 的系数。
\end{exercise}

这些结果通向两个自然方向。Fourier 的 $L^2$ 完备化通向积分方程与 Hilbert 空间；
Riesz--Markov、矩匹配与半群预解式则通向对偶空间、矩问题和算子理论。Banach--Alaoglu
定理、弱紧性、紧算子谱论与谱测度将在后续泛函分析中进一步统一这些方向。
